\section{Mean slopes on deep covers}\label{sec:pf:deepstr}

\providecommand{\deepstrMeanAll}{\operatorname{MeanAll}}
\providecommand{\deepstrMeanGe}{\operatorname{MeanGe}}
\providecommand{\deepstrWtGap}{\operatorname{WtGap}}
\providecommand{\deepstrbnd}{\operatorname{bnd}}
\providecommand{\deepstrerr}{\operatorname{err}}
\providecommand{\deepstrcost}{\operatorname{cost}}
\providecommand{\deepstrclear}{\operatorname{clr}}
\providecommand{\deepstrex}{\operatorname{ex}}
\providecommand{\deepstrspread}{\operatorname{spread}}
\providecommand{\deepstrflagdeg}{\operatorname{fdeg}}
\providecommand{\deepstrdegFl}{\operatorname{deg_{Fl}}}
\providecommand{\deepstrlin}{\operatorname{lin}}
\providecommand{\deepstrslope}{\operatorname{slope}}
\providecommand{\deepstrvol}{\operatorname{vol}}
\providecommand{\deepstrcvol}{\operatorname{cvol}}
\providecommand{\deepstrvsl}{\operatorname{vslope}}
\providecommand{\deepstrtop}{\operatorname{tdeg}}
\providecommand{\deepstrsmall}{\operatorname{Small}}
\providecommand{\deepstrmsmall}{\operatorname{MSmall}}
\providecommand{\deepstrcusp}{\operatorname{Cusp}}
\providecommand{\deepstrComp}{\operatorname{Comp}}
\providecommand{\deepstrlogc}{\operatorname{lc}}
\providecommand{\deepstrII}{\mathrm{II}}
\providecommand{\deepstrarvol}{\widehat{\operatorname{vol}}}
\providecommand{\deepstrgdeg}{\operatorname{gdeg}}
\providecommand{\deepstrrd}{\operatorname{rd}}
\providecommand{\deepstrPrimes}{\mathcal{P}}
\providecommand{\deepstrtr}{\operatorname{tr}}
\providecommand{\deepstrev}{\operatorname{ev}}
\providecommand{\deepstrDim}{D_{\mathrm{Fl}}}

This section proves \cref{thm:deepstr}: the eight published inputs listed there imply the proposition
$\sfname{DeepStR}$ of \cref{def:covolume:deepstr}, a uniform lower bound for the mean, over the connected
components of a deep cover, of arithmetic volume slopes of automorphic line bundles.

\subsection{Introduction}\label{sec:pf:deepstr:intro}

\subsubsection*{What is proved}
The proposition $\sfname{DeepStR}$ is a statement about the unitary moduli spaces $X^o$ of
\cref{sec:covolume}: for every weight $u$ with a gap $g$, the arithmetic volume slopes
$\deepstrvsl(Z,u)$ of the automorphic line bundle of weight $u$ on the components $Z$ of $X^o$ (counting
cusp forms that are small at the archimedean places and integral for prescribed families of valuations at
the finite places) have a mean bounded below by $-C\,\deepstrerr(u,g)$, where the error is linear in the
size $\|u\|_1$ of the weight, logarithmic in the gap, and depends on the host field only through
$\log\deepstrrd_E+|\frac{L'}{L}(1,\chi_d)|$. The constant $C$ is uniform over all global data, core data,
levels and deep covers.

The proof has two parts of very different nature.
\begin{enumerate}[label=(\Alph*)]
\item\emph{Transport from root towers.} The good tower volume theorem (\cref{thm:tower},
\lref{GlobalTheory.CompanionTheory.Ar.Link.good_tower_volume_theorem_bis}) provides, for every level
with good reduction everywhere, a \emph{root tower}: for every bound $B\in\N$ and every component $c$, an
arithmetic variety in multiprojective coordinates (a \emph{root model}) realising a finite cover of a flag
bundle over the component, on which arithmetic intersection numbers of the Chern-root line bundles are
defined; and, for abelian CM fields, it proves a lower bound for the mean of their slopes
(\cref{def:pf:deepstr:towerstatements}). The main step of this section
(\cref{thm:pf:deepstr:comparison}) compares the slope of a root model with the volume slope of the
corresponding component: by an arithmetic Hilbert--Samuel inequality the slope bounds the counting volume
of small sections of the model from below, and these sections are transported, by a trace map, to small
cusp forms on the component; denominators at the finitely many relevant primes are cleared at an explicit
cost $2\deepstrcost(l,B)\|u\|_1$, which is $O((\log B+\log\deepstrrd_E)\|u\|_1)$
(\cref{lem:pf:deepstr:costbound}). Averaging over components and inserting the parameters of deep covers
gives $\sfname{DeepStR}$ (\cref{thm:pf:deepstr:meanall,thm:pf:deepstr:ofvolume}).
\item\emph{Supply of the inputs.} The good tower volume theorem and the comparison need a package of
fourteen arithmetic-geometric statements ($\sfname{Given}$, \cref{def:pf:deepstr:given}), the existence of
lines of large degree, the existence of a full arithmetic intersection system, resolution in projective
coordinates, quotients of abelian schemes by finite flat subgroups and a sharp form of the Hodge--Tate
statement. These are derived from the eight inputs of \cref{thm:deepstr} together with statements proved
in the development.
\end{enumerate}

\subsubsection*{Organisation}
\Cref{sec:pf:deepstr:conv} fixes conventions on counting volumes and recalls the volume slopes of
\cref{sec:covolume}. \Cref{sec:pf:deepstr:towers} defines root models, root towers and their origins, and
states the two tower propositions. \Cref{sec:pf:deepstr:inputs} defines the package $\sfname{Given}$ and
derives all inputs from the eight hypotheses; it ends with the proof of the main theorem.
\Cref{sec:pf:deepstr:volume} reduces $\sfname{DeepStR}$ to a comparison theorem, which is proved in
\cref{sec:pf:deepstr:comparison}. \Cref{sec:pf:deepstr:further} proves further results used here and in other
sections of Part II (symmetry and multilinearity of intersection numbers, intersection numbers of
combinations of lines, nefness on the generic fibre, flag varieties, Hodge--Tate points and charts).
\Cref{sec:pf:deepstr:remarks} collects remarks on the formal statements.

\subsubsection*{Scope of the conversion}
Part (A) is given in full down to the following geometric statements about root models and their origins,
which are cited with their formal statements: the identification of the top degree of a component with
the volume of the root model (\lref{GlobalTheory.T.uniform_norm_degree_bound}), the existence of
generic frames and of a small basis (\lref{GlobalTheory.T.tower_origin_component_bound},
\lref{GlobalTheory.T.generic_frame_component_bound}), the descent of small model sections to cusp forms
(\lref{GlobalTheory.T.uniform_norm_all_tower_bound}), the comparison of the two valuation families
(\lref{GlobalTheory.T.low_height_of_given}), the approximation of the Chern-root metrics by integral
combinations (\lref{GlobalTheory.T.strict_weight_growth_bound}); the identities for intersection numbers
of combinations on root models are proved in \cref{sec:pf:deepstr:further}.
In part (B) the reductions are given in full; the proofs of the derived statements (lines of large degree,
the Lie slope bound of $\sfname{Given}$, the provisions of Hilbert--Samuel type, the Poincar\'e--Lelong formula
in coordinates, quotients by finite flat group schemes, the sharp Hodge--Tate statement) are cited. In
\cref{sec:pf:deepstr:further} the nef polarisation statement and the two Hodge--Tate results are reduced to
cited lemmas on analytic systems, $p$-divisible groups and Hodge components; the other results there are
proved in full, citing besides general facts from the shared part only the symmetry in the last two entries
of field intersection numbers and two integrality statements for chart algebras.

% ---------------------------------------------------------------------------
\subsection{Conventions and volume slopes}\label{sec:pf:deepstr:conv}

\subsubsection*{Counting volumes}
For $m\in\N\cup\{\infty\}$ put $\deepstrlogc(m)=\max(\log m,0)\in[0,\infty]$ (so $\deepstrlogc(0)=0$ and
$\deepstrlogc(\infty)=\infty$). For $D,\delta,a\in\N$ and a sequence $m\colon\N\to\N\cup\{\infty\}$ the
\emph{counting volume} is
\begin{equation}\label{eq:pf:deepstr:countvol}
\deepstrcvol(D,\delta,a;m)=\limsup_{k\to\infty}\frac{D!\,\deepstrlogc(m(ak))}{\delta\,(ak)^{D+1}}
\in[0,\infty],
\end{equation}
and for $\delta>0$ and $v\colon\R\to[0,\infty]$ the \emph{volume slope} is
$\deepstrvsl(\delta,v)=\inf_{t\in\R}\bigl(v(t)/\delta-t\bigr)\in[-\infty,\infty]$. These are the counting volume
and the volume slope used in the proof of \cref{thm:covol} (\lref{GlobalTheory.CompHt.countVol} and
\lref{GlobalTheory.CompHt.volSlope} define them by the two formulas just displayed). For a real number $a$,
\begin{equation}\label{eq:pf:deepstr:volslopeiff}
a\le\deepstrvsl(\delta,v)\iff\text{for all }t\in\R:\ \max\bigl(\delta(a+t),0\bigr)\le v(t);
\end{equation}
this equivalence, immediate from the definition since $\delta>0$, is
\lref{GlobalTheory.CompHt.le_volSlope_iff}.

\subsubsection*{Weights}

\begin{definition}[Weights]\label{def:pf:deepstr:weights}
Let $\Psi$ be a finite set and $N\ge1$. A (real) \emph{weight} is $\lambda\in\R^{\Psi\times\{0,\dots,N-1\}}$; put
$\|\lambda\|=\max_{\tau,i}|\lambda_{\tau,i}|$, $\|\lambda\|_1=\sum_{\tau,i}|\lambda_{\tau,i}|$ and
$\deepstrspread(\lambda)=\sum_\tau(\max_i\lambda_{\tau,i}-\min_i\lambda_{\tau,i})$. The \emph{flag
dimension} is $\deepstrDim=\deepstrDim(\Psi,N)=|\Psi|\,N(N-1)/2$, the dimension of a product over $\Psi$ of
full flag varieties of $\C^N$. Let $\R[X]=\R[X_{\tau,i}]$ be the polynomial ring in the \emph{Chern roots};
put $\deepstrlin(\lambda)=-\sum_{\tau,i}\lambda_{\tau,i}X_{\tau,i}$,
\[
\deepstrdegFl(F)=(-1)^{\deepstrDim}\sum_{\sigma\colon\Psi\to S_N}\Bigl(\prod_\tau\operatorname{sgn}\sigma_\tau\Bigr)
\operatorname{coeff}_{X^{\sigma}}(F)\qquad(F\in\R[X]),
\]
where $X^\sigma=\prod_{\tau,i}X_{\tau,i}^{N-1-\sigma_\tau(i)}$ (this Weyl-alternating coefficient extraction computes
degrees on the flag variety, see \cref{prop:pf:deepstr:flagdeg}), and $\deepstrflagdeg(\lambda)=\deepstrdegFl(\deepstrlin(\lambda)^{\deepstrDim})$;
let $\rho_{\tau,i}=N-1-i$. For $1\le j<N$, $\lambda$ is \emph{regular dominant} if $\lambda_{\tau,b}<\lambda_{\tau,a}$
for $a<b$; its \emph{gap} at $\tau$ is $\lambda_{\tau,j-1}-\lambda_{\tau,j}$; the \emph{small block} is
$\{i<j\}$ if $j\le N-j$ and $\{i\ge j\}$ otherwise. For $g,B,r>0$ and $L\in\R$ put
\begin{equation}\label{eq:pf:deepstr:bnd}
\deepstrbnd_0(\lambda,g)=4|\Psi|N^2\bigl(1+\|\lambda\|/g\bigr),\qquad\Lambda(B,r,L)=\log(Br)+|L|.
\end{equation}
The \emph{weight hypothesis} $\mathrm{WH}(j,\lambda,g)$ means: $g>0$, $\lambda$ is regular dominant, all its
gaps are $\ge g$, and $|\lambda_{\tau,a}-\lambda_{\tau,b}|\ge g$ whenever $a$ is in the small block and $b\ne a$.
\formallabel{GlobalTheory.CompanionTheory.Ar.Top.dimFl}
\formallabel{GlobalTheory.CompanionTheory.Ar.Top.linRoot}
\formallabel{GlobalTheory.CompanionTheory.Ar.Top.degFl}
\formallabel{GlobalTheory.CompanionTheory.Ar.Top.flagDeg}
\formallabel{GlobalTheory.CompanionTheory.Ar.Top.RegDominant}
\formallabel{GlobalTheory.CompanionTheory.Ar.Top.gap}
\formallabel{GlobalTheory.CompanionTheory.Ar.Top.InSmall}
\formallabel{GlobalTheory.CompanionTheory.Ar.Top.spread}
\formallabel{GlobalTheory.CompanionTheory.Ar.Top.bigB0}
\formallabel{GlobalTheory.CompanionTheory.Ar.Top.bigLam}
\formallabel{GlobalTheory.CompanionTheory.Ar.Top.WeightHyp}
\end{definition}

\subsubsection*{Arithmetic varieties in coordinates}
We use the coordinate framework of \cref{def:covolume:hs}: a finite family of blocks of variables
$x_{k,i}$, the multigraded ring $R=\Z[x_{k,i}]$, a point $y$ of a field $\Omega$ with coordinates $y_{k,i}$,
its ideal $P=P_y$ (kernel of $x_{k,i}\mapsto y_{k,i}T_k$), Hermitian lines $A=(a,h)$ (a multidegree and a
positive homogeneous function on the complex multicone), and intersection systems $\deepstrII$. We write
$\deepstrII(P;A^{D+1})$ for the self-intersection of $A$ on an arithmetic variety $P$ of dimension $D+1$
and $\deepstrgdeg(P,A)$ for the geometric degree of $A$ on its generic fibre. For a multidegree $a$ and
$\tau\in\N$, the \emph{chart lattice} $\Lambda_\tau(y,a)\subseteq\Omega$ is the group of $s\in\Omega$ such that
$s\prod_ky_{k,c_k}^{-\tau a_k}$ is integral over the chart ring $\Z[y_{k,i}/y_{k,c_k}]$ for every chart
$c$ with all $y_{k,c_k}\ne0$. An element of $\Lambda_\tau(y,a)$ is \emph{small of degree $\tau$} for $A=(a,h)$ if
its supremum norm for the metric $h^\tau$ on the complex points is at most $1$; the \emph{arithmetic volume}
$\deepstrarvol(A)$ is the $\limsup_\tau$ of the logarithm of the number of small elements of degree $\tau$
divided by $\tau^{D+1}/(D+1)!$. For $t\in\R$ the \emph{shift} $A(t)$ is $A$ with $h$ replaced by $e^th$.
A Hermitian line is \emph{nef} on $P$ if it is semipositive relative to some ideal $P_0\subseteq P$, nef on
the fibres and conjugation-invariant.

\subsubsection*{Unitary moduli spaces and volume slopes}
We recall, in the form needed here, the objects of \cref{def:covolume:deepstr}; their formal definitions
belong to the proof of \cref{thm:covol}. Let $E$ be a CM field with a host structure $H$
(\cref{def:host}), $\Psi=\Psi_H$ the CM type of $H$, $D_w$ a datum of signature $(p,q)$, $N=p+q$, $l$ a
level with level modulus $M'$ and $A$-part $l_A$, $X^o$ a unitary moduli space, $c$ a coefficient datum
with field $F_c\subset\C$ (finite over $\Q$), and $\deepstrComp=\deepstrComp(X^o,c)$ the set of irreducible
components of the generic fibre $X^o_{F_c}$. A family $U=(U_\sigma)_{\sigma\colon F_c\to\C}$ of complex
uniformisations, with base points $\xi_0$, provides frames on each component and Hodge-type norms
$\|f\|_{G,\sigma,z}$ of rational sections $f$ at frames $z$, for weights $G$ of automorphic forms; to an
integer weight $u\in\Z^{\Psi\times\{0,\dots,N-1\}}$ is attached an automorphic weight $G_u$, with
$G_{nu}=nG_u$ for $n\in\N$. For a component $Z$, $\deepstrcusp(Z,G)$ is the $F_c$-space of cusp forms of weight
$G$ supported on $Z$.

The development attaches to each height-one prime $\lambda$ of $\OO_{F_c}$ two sets $\mathcal V^{\mathrm{ht}}(\lambda)$,
$\mathcal V^{\mathrm{rec}}(\lambda)$ of extended valuations $w$ on rational sections ($w(0)=+\infty$,
$w(f+f')\ge\min(w(f),w(f'))$, $w(ff')=w(f)+w(f')$, $w(af)=v_\lambda(a)+w(f)$ for $a\in F_c^\times$); the
first comes from Hodge--Tate frames, the second from reductions. For $B>0$ let
\begin{equation}\label{eq:pf:deepstr:htprimes}
\deepstrPrimes(l,B)=\{p\text{ prime}:\ p\le B\ \text{or}\ p\mid2|d_E|M'\},
\end{equation}
and let $\mathcal V_B(\lambda)$ be $\mathcal V^{\mathrm{ht}}(\lambda)$ if $\lambda$ lies over a prime of
$\deepstrPrimes(l,B)$ and $\mathcal V^{\mathrm{rec}}(\lambda)$ otherwise. For $R\in\R$ put
\begin{multline*}
\deepstrsmall_B(Z,G,R)=\bigl\{f\in\deepstrcusp(Z,G):\ w(f)\ge0\text{ for all }\lambda\text{ and }w\in\mathcal V_B(\lambda),\\
\|f\|_{G,\sigma,z}\le R\ \text{for all good }\sigma\text{ and frames }z\text{ on }Z\bigr\},
\end{multline*}
the \emph{small cusp forms}. Write
\[
\mathrm{vdm}(u)=\prod_\tau\prod_{i<i'}(u_{\tau,i}-u_{\tau,i'}),\qquad
\mathrm{gr}(Z,G)=\limsup_{\tau\to\infty}\frac{\deepstrDim!\,\dim_{F_c}\deepstrcusp(Z,\tau G)}{\tau^{\deepstrDim}}.
\]
The \emph{top degree} is $\deepstrtop(Z,u)=v_0(Z)\,\mathrm{vdm}(u)$, where
$v_0(Z)=\mathrm{gr}(Z,G_{u_0})/\mathrm{vdm}(u_0)$ for a fixed standard weight $u_0$. With
\[
\mathrm{cv}_{Z,u}(t)=\deepstrcvol\bigl(\deepstrDim,[F_c:\Q],8;\ \tau\mapsto\#\deepstrsmall_B(Z,\tau G_u,e^{\tau t})\bigr),
\]
the \emph{volume slope} of $Z$ at $u$ is
\begin{equation}\label{eq:pf:deepstr:compvol}
\deepstrvsl(Z,u)=\deepstrvsl\bigl(\deepstrtop(Z,u),\mathrm{cv}_{Z,u}\bigr)
\end{equation}
(\lref{GlobalTheory.CompHt.compVolSlope} defines it by this formula). With
\begin{gather*}
\deepstrbnd(u,g)=4|\Psi|N^2\bigl(1+\max_{\tau,i}|u_{\tau,i}|/g\bigr),\qquad
\deepstrerr_s(u,g)=(\log\deepstrbnd(u,g)+s)\|u\|_1+q\,\deepstrspread(u),
\end{gather*}
$\deepstrWtGap(u,g)$ means $g>0$, every $u_\tau$ strictly decreasing, and $|u_{\tau,a}-u_{\tau,b}|\ge g$
whenever $a\ge p$ and $b\ne a$; $\deepstrMeanGe(u,B,X)$ means that there are real numbers
$\mathrm{sl}(Z)\le\deepstrvsl(Z,u)$, $Z\in\deepstrComp$ (computed with $\mathcal V_B$), with
$\frac1{\#\deepstrComp}\sum_Z\mathrm{sl}(Z)\ge-X$; and $\deepstrMeanAll(s,C)$ means
$\deepstrMeanGe(u,\deepstrbnd(u,g),C\,\deepstrerr_s(u,g))$ for all $(u,g)$ with $\deepstrWtGap(u,g)$
(\lref{GlobalTheory.CovHt.MeanAll} is this definition). $\sfname{DeepStR}$ is then the statement of
\cref{def:covolume:deepstr}, with $\ell(\mathcal D)=\log\deepstrrd_E+|\frac{L'}{L}(1,\chi_d)|$
(\lref{GlobalTheory.CompanionTheory.Ar.Plug.DeepStR}).

\begin{definition}[Conditions on uniformisations]\label{def:pf:deepstr:unifconditions}
Let $(X^o,c,U)$ be as above.
\begin{enumerate}[label=(\alph*)]
\item $U$ satisfies the \emph{composition laws} (CL) with respect to a choice $h$ of generic frames of the
components and of sections $e$ if, among other conditions, the number of components of the uniformisation at
the identity embedding equals the number of components of $X^o_{F_c}$, $X^o_{F_c}$ has a component, the generic
frames are obtained from test frames, evaluation of sections at frames is well defined, and the $U_\sigma$ for
good $\sigma$ are uniformisations by test frames.
\item $(X^o,c)$ has the \emph{uniformisation property} (UP) if every algebraic integral frame over a
complete discrete valuation ring above a height-one prime of $\OO_{F_c}$ becomes, after a finite extension of
the ring, an algebraic frame which admits a Mumford-type degeneration datum with a positive semi-definite
monodromy form.
\item $(X^o,c)$ is \emph{geometrically irreducible} (GI) if $F_c$ is algebraically closed in the function
field of every component.
\end{enumerate}
\formallabel{GlobalTheory.CompanionTheory.Ar.Rec.CompLaws}
\formallabel{GlobalTheory.CompanionTheory.Ar.Ref.UnifAll}
\end{definition}

% ---------------------------------------------------------------------------
\subsection{Root models, root towers and the tower propositions}\label{sec:pf:deepstr:towers}

\begin{definition}[Root models]\label{def:pf:deepstr:rootmodel}
Let $\Psi$ be a finite set, $N\ge1$, $\deepstrDim=\deepstrDim(\Psi,N)$, and let $\mathcal R=(\Psi\times\{0,\dots,N-1\})\sqcup\{\ast\}$.
A \emph{root model} $M$ of type $(\Psi,N)$ consists of
\begin{enumerate}[label=(\roman*)]
\item coordinates (finite sets $K$ and $\iota_k$, $k\in K$), a field $\Omega_M$ and a point $y\in\Omega_M^{\sqcup_k\iota_k}$
whose ideal $P_M=P_y$ is an arithmetic variety of dimension $\deepstrDim+1$, $y$ being \emph{birational}
(every element of $\Omega_M$ is a quotient of two multihomogeneous polynomials of the same multidegree
evaluated at $y$);
\item multidegrees $a_r\in\Z^K$ and, for every $t\in\R$, smooth conjugation-invariant Hermitian lines
$x_t(r)$ of multidegree $a_r$, for $r\in\Psi\times\{0,\dots,N-1\}$ (the \emph{Chern-root lines});
\item a positive integer $k_A$ and a smooth conjugation-invariant Hermitian line $A$; we put
$x_t(\ast)=A$;
\item a positive integer $\delta_M$.
\end{enumerate}
For $n\in\Z^{\mathcal R}$ let $L_t(n)=\bigotimes_rx_t(r)^{\otimes n_r}$ (multidegree $\sum_rn_ra_r$, metric
$\prod_rh_{x_t(r)}^{n_r}$). For a weight $\lambda$ and $s\in\R$ let $\mathrm{coef}(\lambda,s)\in\R^{\mathcal R}$ be
$-\lambda_r$ at $r\ne\ast$ and $s/k_A$ at $\ast$. On the polynomial ring $\R[X_r:r\in\mathcal R]$, an
intersection system $\deepstrII$ and $t\in\R$ define linear functionals $\mathrm{ar}_t$ and $\mathrm{geo}$
by their values on monomials $X^m$: $\mathrm{ar}_t(X^m)=\deepstrII(P_M;(x_t(r))^{m})/(\delta_Mk_A^{m_\ast})$
if $\deg m=\deepstrDim+1$ (each line repeated $m_r$ times), and $\mathrm{geo}(X^m)$ is the integral over the
complex points of $P_M$ of the product of the Chern forms of the lines $x_1(r)$, repeated $m_r$ times,
divided by $\delta_Mk_A^{m_\ast}$, if $\deg m=\deepstrDim$; both vanish on monomials of other degrees.
\formallabel{GlobalTheory.CompanionTheory.Ar.Top.RootModel}
\formallabel{GlobalTheory.CompanionTheory.Ar.Top.RootModel.geoFun}
\end{definition}

The \emph{limit functional} $\mathrm{PP}=\mathrm{PP}_M(\deepstrII)$ is defined on monomials by
$\mathrm{PP}(X^m)=\lim_{t\to0^+}\mathrm{ar}_t(X^m)$ and extended linearly
(\lref{GlobalTheory.CompanionTheory.Ar.Top.RootModel.PP} is this definition, the value being an
unspecified real number when the limit does not exist).

\begin{definition}[Root towers]\label{def:pf:deepstr:roottower}
A \emph{root tower} $T$ of type $(\Psi,N)$ consists of an integer $j$ with $1\le j<N$, a finite nonempty
set $\deepstrComp(T)$ of \emph{components}, a finite set $\mathrm{bad}(T)\subset\N$ containing $2$, a
rational weight $\nu$, and a root model $M(B,c)$ of type $(\Psi,N)$ for every $B\in\N$ and
$c\in\deepstrComp(T)$. Put $S_T(B)=\{p\text{ prime}:\ p\le B\text{ or }p\in\mathrm{bad}(T)\}$, and, for an
intersection system $\deepstrII$, a weight $\lambda$ and $s>0$ (all functionals being those of $M(B,c)$),
\begin{gather*}
\deepstrvol_T(c)=\frac{\mathrm{geo}_{M(2,c)}(\deepstrlin(\rho)^{\deepstrDim})}{\deepstrflagdeg(\rho)},\qquad
\deepstrslope_T(B,c,\lambda)=\frac{\mathrm{PP}_{M(B,c)}(\deepstrlin(\lambda)^{\deepstrDim+1})}
{(\deepstrDim+1)\,\deepstrvol_T(c)\,\deepstrflagdeg(\lambda)},\\
\lambda_s=\lambda-s\nu,\qquad \ell_s(\lambda)=\deepstrlin(\lambda_s)+sX_\ast,\\
\overline m(B,c,\lambda,s)=\mathrm{ar}_{\varsigma(s)}\bigl(\ell_s(\lambda)^{\deepstrDim+1}\bigr),\qquad
m(B,c,\lambda,s)=\mathrm{geo}\bigl(\ell_s(\lambda)^{\deepstrDim}\bigr),
\end{gather*}
with $\varsigma(s)=\exp(-s^{-2})$, and for $r,L\in\R$
\[
\beta_T(r,L,B,\lambda)=\Lambda(B,r,L)\,\|\lambda\|_1+\min(j,N-j)\,\deepstrspread(\lambda).
\]
We use the following conditions on $T$ and a full intersection system $\deepstrII$.
\begin{itemize}
\item \emph{Proportionality}: $\deepstrvol_T(c)>0$, $\mathrm{geo}_{M(B,c)}(F)=\deepstrvol_T(c)\deepstrdegFl(F)$ for
all $F\in\R[X]$, and $\mathrm{geo}_{M(B,c)}$ does not depend on $B$.
\item \emph{Equality of volumes}: $\deepstrvol_T(c)$ does not depend on $c$.
\item The \emph{limit conditions} $\mathrm{(LC)}$: for all $(\lambda,g)$ with $\mathrm{WH}(j,\lambda,g)$, all
$B\in\N$ with $\deepstrbnd_0(\lambda,g)<B+1$ and all $c$: for all sufficiently small $s>0$ the model $M(B,c)$
is nef on the fibres and semipositive at $(\lambda_s,s)$; as $s\to0^+$,
$\overline m(B,c,\lambda,s)\to\mathrm{PP}(\deepstrlin(\lambda)^{\deepstrDim+1})$,
$m(B,c,\lambda,s)\to\deepstrvol_T(c)\deepstrflagdeg(\lambda)$ and
$\overline m/((\deepstrDim+1)m)\to\deepstrslope_T(B,c,\lambda)$; and
$\deepstrvol_T(c)\deepstrflagdeg(\lambda)>0$.
\item The \emph{mean-slope bound} $\mathrm{(MS)}$ with constants $(C,r,L)$: for all $(\lambda,g)$ with
$\mathrm{WH}(j,\lambda,g)$ and all $B\in\N$ with $B\ge128$, $\deepstrbnd_0(\lambda,g)<B+1$ and
$\#\mathrm{bad}(T)\le B$,
\begin{equation}\label{eq:pf:deepstr:meanslope}
\frac1{\#\deepstrComp(T)}\sum_{c}\deepstrslope_T(B,c,\lambda)\ \ge\ -C\,\beta_T(r,L,B,\lambda).
\end{equation}
\end{itemize}
\formallabel{GlobalTheory.CompanionTheory.Ar.Top.RootTower}
\formallabel{GlobalTheory.CompanionTheory.Ar.Top.RootTower.InS}
\formallabel{GlobalTheory.CompanionTheory.Ar.Top.RootTower.bracket}
\formal{GlobalTheory.CompanionTheory.Ar.Top.RootTower.T1}
\formal{GlobalTheory.CompanionTheory.Ar.Top.RootTower.MeanSlope}
\end{definition}

The volume, the slope and proportionality are defined in \lref{GlobalTheory.CompanionTheory.Ar.Top.RootTower.vol},
\lref{GlobalTheory.CompanionTheory.Ar.Top.RootTower.slope} and
\lref{GlobalTheory.CompanionTheory.Ar.Top.RootTower.Proportional} by the formulas above.

\begin{definition}[Origins, small model sections and readings]\label{def:pf:deepstr:origins}
Let $(E,H,D_w,l,X^o,c)$ be as in \cref{sec:pf:deepstr:conv} and let $Y$ be a model of $X^o$ over $\OO_{F_c}$
(a normal scheme, projective and flat over $\OO_{F_c}$, with an identification of its generic fibre with
$X^o_{F_c}$). Let $T$ be a root tower of type $(\Psi_H,p+q)$.
\begin{enumerate}[label=(\alph*)]
\item A \emph{root origin} $\mathrm{Or}$ of a root model $M$ over $Y$, relative to a finite set $S$ of
primes and to fixed data $F_x$ on $Y$, consists of a Hodge component of $Y$ at $S$ with a flag $\mathrm{Fl}$
of degree $\deg\mathrm{Fl}\ge0$; a field embedding of the function field $\Omega_f$ of $F_x$ into $\Omega_M$ with
$[\Omega_M:\Omega_f]=\deg\mathrm{Fl}$, compatible with the flag; the identity
\begin{equation}\label{eq:pf:deepstr:deltaeq}
\delta_M=[F_c:\Q]\cdot\deg\mathrm{Fl};
\end{equation}
Hermitian lines $\mathrm{sing}(r)$ of multidegree $a_r$ which are the Hodge lines of the flag; the
identification of $A$ with the pull-back of the Hodge-type line of $F_x$; an effectivity condition; and a
realisation of the analytic system of $F_x$ on the complex points of $P_M$, relating the lines
$\mathrm{sing}(r)$ and $x_t(r)$ through fundamental weights.
\item A \emph{tower origin} $o$ of $T$ on $Y$ consists of: the equality $j=p$; the equivalence
$r\in\mathrm{bad}(T)\iff r$ prime and $r\mid2d_EM'$; the sets $S(B)=S_T(B)$; for each $a\in\deepstrComp(T)$
fixed data $F_x(a)$ on $Y$, whose generic points are pairwise distinct and exhaust the generic points of
$Y$; a compatibility of $\nu$ with $F_x(a)$; and root origins $\mathrm{Or}(B,a)$ of $M(B,a)$ relative to
$S(B)$ and $F_x(a)$. A bijection $\kappa\colon\deepstrComp(X^o,c)\to\deepstrComp(T)$ is \emph{compatible} with $o$ if,
for every component $Z$, the generic point of $F_x(\kappa(Z))$ is the image in $Y$ of the generic point of $Z$.
\item A \emph{good origin} is a tower origin whose Hodge components carry vector bundles, whose
realisation is smooth and flat, together with reference root models $R(a)$, root origins of degree $1$
over $S=\varnothing$ and with an ample line $A$. The \emph{immersion condition} $\mathrm{(IC)}$ on a good
origin is a local immersion condition on the analytic systems together with a numerical condition at the
primes of $S(B)$.
\item For an integer weight $u$ let $n_u\in\Z^{\mathcal R}$ be $n_u(r)=-u_r$ for $r\ne\ast$ and $n_u(\ast)=0$,
and $a_u=\sum_rn_u(r)a_r$. The \emph{small model sections} are
$\deepstrmsmall(\mathrm{Or},u,R)$ is the set of $s\in\Lambda_1(y,a_u)$ such that, for some multihomogeneous form
$w\notin P_M$,
\[
|s|(z)\le R\prod_rh_{\mathrm{sing}(r)}(z)^{n_u(r)}\qquad\text{at all complex points $z$ of $P_M$ with }w(z)\ne0,
\]
$|s|(z)$ being the value of $s$ as a section of multidegree $a_u$.
\item Let $Z$ be a component of $X^o_{F_c}$ and $\Omega(Z)$ its function field; for a rational section $f$ let
$\deepstrev(f)\in\Omega(Z)$ be its value at the generic frame of $Z$. A \emph{reading} of a generic frame of
$\mathrm{Or}$ on $Z$ consists of a field embedding $\varphi\colon\Omega(Z)\to\Omega_M$, compatible with constants, with
$[\Omega_M:\varphi\Omega(Z)]=\deg\mathrm{Fl}$, and units of $\Omega_M$ relating the frame to $\varphi\circ\deepstrev$. For an
integer weight $v$ it defines the \emph{trace} $\deepstrtr_v(s)=\Tr_{\Omega_M/\Omega(Z)}(s/\gamma_v)$, $s\in\Omega_M$, with
a gauge $\gamma_v\in\Omega_M^\times$.
\end{enumerate}
\formallabel{GlobalTheory.CompanionTheory.Ar.RootOrigin}
\formallabel{GlobalTheory.CompanionTheory.Ar.TowerOrigin}
\formallabel{GlobalTheory.CompanionTheory.Ar.Top.GoodOrigin}
\formallabel{GlobalTheory.CompanionTheory.Ar.Plug.degOf}
\formallabel{GlobalTheory.CompanionTheory.Ar.Plug.HodgeLe}
\formallabel{GlobalTheory.CompanionTheory.Ar.Plug.modelSmall}
\formallabel{GlobalTheory.CompanionTheory.Ar.Rec.Reading}
\formallabel{GlobalTheory.CompanionTheory.Ar.Plug.trOf}
\end{definition}

\begin{definition}[The tower propositions]\label{def:pf:deepstr:towerstatements}
For a host $H$ of $E$ let $\deepstrrd_E=|d_E|^{1/[E:\Q]}$ and $L_H=\Re\frac{L'(1,\chi_d)}{L(1,\chi_d)}$,
$\chi_d=(\frac{\cdot}{d})$. A model $Y$ of $X^o_l$ is \emph{global} if $l$ has good reduction at all primes,
$l_A=1$, and $Y$ descends and dominates the local models provided by the toroidal package
(\cref{def:pf:deepstr:given}\ref{G:pf:deepstr:tor}) at every prime.
\begin{enumerate}[label=(\alph*)]
\item $\sfname{GoodTowerExistsX}$: for every CM field $E$, host $H$, datum $D_w$ with $p,q\ge1$, family
$(X^o_{l'})_{l'}$ of moduli spaces indexed by the levels, level $l$, coefficient datum $c$ and global model
$Y$ of $X^o_l$ over $\OO_{F_c}$ with split components, there are a root tower $T$ of type $(\Psi_H,p+q)$ and a
good origin $o$ of $T$ on $Y$ satisfying $\mathrm{(IC)}$.
\item $\sfname{VolumeTheoremX}$: there is $C\ge0$ such that for every \emph{abelian} CM field $E$, every
$H,D_w,(X^o_{l'}),l,c$, global model $Y$, root tower $T$ with a good origin $o$ satisfying $\mathrm{(IC)}$, and
full intersection system $\deepstrII$: if $q\ge1$, $3q+1\le p+q$ and $\#\mathrm{bad}(T)\le4|\Psi_H|(p+q)^2$,
then $T$ satisfies proportionality, equality of volumes, $\mathrm{(LC)}$, an asymptotic formula for the limit
functionals, a lower bound for their sum over the components, and $\mathrm{(MS)}$ with constants
$(C,\deepstrrd_E,L_H)$.
\end{enumerate}
\formallabel{GlobalTheory.CompanionTheory.Ar.Top.GoodTowerExistsX}
\formallabel{GlobalTheory.CompanionTheory.Ar.Top.VolumeTheoremX}
\end{definition}

The tower propositions are the content of \cref{thm:tower}: under the hypotheses listed there,
$\sfname{GoodTowerExistsX}\wedge\sfname{VolumeTheoremX}$ holds
(\lref{GlobalTheory.CompanionTheory.Ar.Link.good_tower_volume_theorem_bis}). Of
$\sfname{VolumeTheoremX}$ this section uses only proportionality, $\mathrm{(LC)}$ and $\mathrm{(MS)}$.

% ---------------------------------------------------------------------------
\subsection{The inputs and the proof of the main theorem}\label{sec:pf:deepstr:inputs}

\begin{hypothesis}[The eight inputs]\label{hyp:pf:deepstr:inputs}
The following statements are those of \cref{thm:deepstr}, in the form of the development; schemes are
over a field $k$ of characteristic $0$ in (H1), (H2).
\begin{enumerate}[label=(H\arabic*)]
\item\label{H:pf:deepstr:res} \emph{Resolution}: every reduced scheme $X$, locally of finite type,
quasi-compact and separated over $k$, admits a sequence of blow-ups $s\colon X'\to X$ with $X'$ smooth over
$k$, $s$ an isomorphism over the smooth locus of $X$, and the reduced preimage of the non-smooth locus a
divisor with strict normal crossings \cite{Hironaka1964,BierstoneMilman1997,Wlodarczyk2005,Kollar2007}.
\item\label{H:pf:deepstr:princ} \emph{Principalisation}: for $X$ smooth of pure relative dimension $d$,
quasi-compact and separated over $k$, an ideal sheaf $I$ whose support contains no irreducible component of
$X$, and a closed subset $F$ with strict normal crossings, there is a sequence of blow-ups $s\colon X'\to X$
with centres having strict normal crossings with the boundary, $X'$ smooth, $I\OO_{X'}$ an effective Cartier
divisor with strict normal crossings support, and $s$ an isomorphism over the complement of the support of
$I$ (same references).
\item\label{H:pf:deepstr:pl} \emph{Poincar\'e--Lelong}: an integrated form of the arithmetic
Poincar\'e--Lelong formula in coordinates: for an arithmetic variety $P$ of dimension $n+2$, smooth
conjugation-invariant Hermitian lines $L_1,\dots,L_n,A,B$ with $B$ of multidegree $0$, and a section $s$ of
$A$ on $P$, the integral of $\log\|s\|_A$ against $dd^c\log h_B\wedge c_1(L_1)\wedge\dots\wedge c_1(L_n)$ over
the complex points of $P$ equals $\sum_W\mathrm{mult}_W(P+(s))\int_{W(\C)}\log h_B\,c_1(L_1)\cdots c_1(L_n)$,
the sum over the horizontal components $W$ of $P+(s)$, minus $\int_{P(\C)}\log h_B\,c_1(A)c_1(L_1)\cdots c_1(L_n)$
\cite{GilletSoule1990,BostGilletSoule1994}.
\item\label{H:pf:deepstr:hs} $\sfname{HSUse}$ (\cref{def:covolume:hs}(d)).
\item\label{H:pf:deepstr:mor} $\sfname{MoriwakiBigStatement}$ (\cref{def:covolume:hs}(d)) \cite{Moriwaki2009,Yuan2008}.
\item\label{H:pf:deepstr:far} \emph{Exactness of the Hodge--Tate sequence}: for a prime $p$, an
algebraically closed field $C$ of characteristic $0$ with a valuation $v$ (written multiplicatively),
$v(p)<1$, whose valuation ring $\OO$ is $p$-adically complete, a $p$-divisible group $G$ over $\OO$, a
compatible system $(\zeta_n)$ of $p$-power roots of unity with $\zeta_1\ne1$, and the $\OO$-bilinear form on
$(T_pG\otimes\OO)\times(T_pG^D\otimes\OO)$ induced by the Weil pairing: an element lies in the kernel of the
Hodge--Tate map $\alpha_G\otimes\OO$ if and only if it is orthogonal to the kernel of the dual map
$\mathrm{dlog}$ \cite{Faltings1988,FarguesGenestierLafforgue2008}.
\item\label{H:pf:deepstr:zar} \emph{Zarhin's trick}: over the valuation ring $R$ of an algebraically closed
valued field, every abelian scheme with a relatively ample line bundle is a retract, by homomorphisms
$i\colon A\to A'$, $r\colon A'\to A$ with $r\circ i=\mathrm{id}$, of a principally polarised abelian scheme $A'$
\cite{Zarhin1975,Zarhin1985}.
\item\label{H:pf:deepstr:nef} \emph{Asymptotic Riemann--Roch for nef divisors}: for an arithmetic variety
$P=P_y$ of dimension $D+1$ with $y$ birational over a field of characteristic $0$, and a multidegree $a$ nef on
the generic fibre, the $\Q$-spans $V_\tau$ of $\Lambda_\tau(y,a)$ are finite-dimensional and
$D!\dim_\Q V_\tau/\tau^D$ tends to the geometric degree of $a$ \cite{Lazarsfeld2004}.
\end{enumerate}
\formal{GlobalTheory.CompanionTheory.Hypotheses₂.ResolutionCharZero}
\formal{GlobalTheory.CompanionTheory.Hypotheses₂.PrincipalisationCharZero}
\formal{GlobalTheory.CompanionTheory.Hypotheses₂.PoincareLelongInv}
\formal{GlobalTheory.CompanionTheory.Ar.HT.FarguesExact}
\formal{GlobalTheory.CompanionTheory.Ar.GR.ZarhinStatement}
\formal{GlobalTheory.CompanionTheory.Ar.Pub.NefRR}
\end{hypothesis}

\begin{theorem}[\cref{thm:deepstr}]\label{pf:deepstr:main}
Assume \ref{H:pf:deepstr:res}--\ref{H:pf:deepstr:nef}. Then $\sfname{DeepStR}$ holds: there is $C\ge0$ such
that for every global datum $\mathcal D$, every core datum $T$ of $\mathcal D$ satisfying $\mathsf{LawsJ}$ with
$m^2\le k$, every $1\le j<m$ and every deep cover $Y$ of $T$ at level $j$ which is good for some parameters
and whose level modulus has at most $15$ prime factors, $\deepstrMeanAll(\ell(\mathcal D),C)$ holds on $Y$.
\formallabel{PositivityTheory.HtPt.Ah.deep_slim_use}
\end{theorem}

The proof is given at the end of this subsection, after the derivation of the intermediate inputs.

\subsubsection*{The package $\sfname{Given}$}

\begin{definition}[$\sfname{Given}$]\label{def:pf:deepstr:given}
$\sfname{Given}$ is the conjunction of the following fourteen statements.
\begin{enumerate}[label=(G\arabic*)]
\item\label{G:pf:deepstr:tor} \emph{The toroidal package}: a moduli space exists for every level; families
of complex uniformisations of every admissible type exist; for every prime $p$ and valuation subring
$W\subset\C$ over $p$ there are local systems of integral models with split fibres, and global models
dominate them; for every host and datum, change-of-level maps are finite \'etale and surjective on generic
fibres with fibres the orbits of the level group (\emph{level compatibility}), and points have unique
Hecke translates by subgroups of the $p$-torsion at split primes $p\nmid M'$ (\emph{Hecke compatibility}).
\item\label{G:pf:deepstr:key} \emph{The theta-constant formula}: the existence of a function on principally
polarised abelian varieties with theta structure, rigidification and Lie basis, with a transformation law
under change of theta function and Lie basis, compatible with base change, integral over discrete valuation
rings, and normalised by the classical theta functions with characteristics over $\C$.
\item\label{G:pf:deepstr:sah} \emph{Extension of homomorphisms} of semi-abelian schemes over a dense open
subset of a normal noetherian base \cite[Ch.~I, Prop.~2.7]{FaltingsChai1990}.
\item\label{G:pf:deepstr:segre} \emph{The Segre formula}: a formula for intersection numbers on projective
bundles in coordinates in terms of Segre-type sums over the base.
\item\label{G:pf:deepstr:resk} \emph{Resolution in projective coordinates} $\mathrm{(RPC)}$
(\cref{def:pf:deepstr:resk}).
\item\label{G:pf:deepstr:sa} \emph{Strong approximation} for special unitary groups: if $E/F$ is a
quadratic extension of number fields and $J\in\GL_N(E)$ is hermitian and isotropic, then
$\mathrm{SU}(J)(\A_{F,f})$ is contained in the closure of $\mathrm{SU}(J)(F)$ \cite{Kneser1966,PlatonovRapinchuk1994}.
\item\label{G:pf:deepstr:ht} \emph{The sharp Hodge--Tate statement}: the conclusion of \ref{H:pf:deepstr:far},
and in addition $\kappa\cdot\omega_{G^D}\subseteq\operatorname{im}(\alpha_G\otimes\OO)$ for every $\kappa\in\OO$ with
$v(\kappa)^{p-1}\le v(p)$ if $p$ is odd and $v(\kappa)\le v(4)$ if $p=2$ \cite{Faltings1988,FarguesGenestierLafforgue2008}.
\item\label{G:pf:deepstr:gaga} \emph{GAGA for sections}: on a smooth proper complex scheme, every family of
local holomorphic functions compatible with the transition functions of a line bundle is an algebraic
section \cite{Serre1956}.
\item\label{G:pf:deepstr:slope} \emph{The Lie slope bound}: there is $C\ge0$ such that for every CM field $E$,
host $H$, number field $L$ and semi-abelian scheme $t$ over $\OO_L$ with $H$-action and principally polarised
generic fibre, every Hermitian bundle $\overline{\mathcal B}$ over $L$ which is a fibre of the Hodge inner product
on the Lie algebra of $t$ satisfies $\mu_{\max}(\overline{\mathcal B})\le
C\log\bigl(2+\rank\mathcal B+\max(-\widehat{\deg}_n\overline{\mathcal B},0)\bigr)$, $\widehat{\deg}_n$ the normalised
Arakelov degree (compare \cite{Bost1996}).
\item\label{G:pf:deepstr:pb} \emph{Potential-theoretic inputs}: fifteen statements on Green potentials in
coordinates (convergence, signed convergence and its rate, uniform integrability, pull-back, global
regularisation), on resolution and degrees, on non-negativity, rationality and locality of arithmetic
intersection numbers, on the projection formula, fibre lines and fibre dimensions, two Kleiman-type
statements and a Siu-type statement for nef lines.
\item\label{G:pf:deepstr:provhs} \emph{Hilbert--Samuel provisions}: nine statements: the lower and upper
successive-minima bounds of Zhang \cite{Zhang1995}, non-negativity of arithmetic self-intersections, a lower
bound of arithmetic volumes by self-intersections, the volume inequality $\mathrm{(VI)}$ of
\cref{def:pf:deepstr:bppf}, two arithmetic Hilbert--Samuel upper bounds, a comparison of the
self-intersections of two lines related by injections of section lattices with dominated norms, and a
variant of this comparison involving arithmetic volumes.
\item\label{G:pf:deepstr:provan} \emph{Analytic provisions}: seven statements on Green potentials (tail bounds,
convergence, signed convergence, convergence rate, uniform integrability, pull-back, global regularisation).
\item\label{G:pf:deepstr:zhang} \emph{Zhang-type provisions}: the conjunction of Zhang's lower bound, a
dimension statement, the true proposition, a push-down statement and a Grassmannian height bound.
\item\label{G:pf:deepstr:anti} \emph{Antisymmetric perturbation}: if a real polynomial $h$ in the variables
$X_{\tau,i}$ vanishes on compact sets $K_\tau$ contained in the hyperplanes $X_{\tau,j-1}=X_{\tau,j}$ and
containing nonempty relatively open subsets of them, there is $C$ such that every polynomial $P$ with
$P-vh$ antisymmetric and $|P|\le\eps$ on all $K_\tau$ differs from $vh+\Delta\cdot l$, for some invariant
$l$ ($\Delta$ the Vandermonde product), by a polynomial whose coefficients are at most $C\eps$.
\end{enumerate}
\formallabel{GlobalTheory.CompanionTheory.Ar.Top.Given}
\formallabel{GroundTheory.Heights₂.Heights.NumbersDVR.NumericalSegre.Statement}
\formallabel{ReciprocityTheory.ChartRestriction.G.SA.SAStatement}
\formallabel{GlobalTheory.CompanionTheory.Ar.HT.FarguesHodgeTate}
\formallabel{GlobalTheory.CompanionTheory.Ar.SlopeBoundSt}
\formallabel{GlobalTheory.CompanionTheory.Ar.Prov.HS}
\formallabel{GlobalTheory.CompanionTheory.Ar.Prov.An}
\formallabel{GlobalTheory.CompanionTheory.Ar.Prov.ZhangDimensionPushDownGrassmannianHeightBounds}
\formallabel{GroundTheory.Heights₂.Heights.AntisymPerturbation.Statement}
\end{definition}

The components \ref{G:pf:deepstr:tor}, \ref{G:pf:deepstr:key} and \ref{G:pf:deepstr:pb} are the shared
statements \lref{ReciprocityTheory.ChartRestriction.G.Geo.Tor.PartI} (with the level, Hecke and
uniformisation statements next to it), \lref{ReciprocityTheory.ChartRestriction.G.Geo.Tor.KeyFormula}
and \lref{GlobalTheory.CompanionTheory.Ar.PB.Inputs}, whose precise formulations we have summarised.

\begin{definition}[The volume inequality]\label{def:pf:deepstr:bppf}
The \emph{volume inequality} $\mathrm{(VI)}$ is the statement: for every arithmetic variety $P=P_y$ of
dimension $D+1\ge2$ with $y$ birational, every intersection system $\deepstrII$ and every conjugation-invariant
Hermitian line $A$ which is nef on the fibres and semipositive relative to some $P_0\subseteq P$,
\[
\deepstrII(P;A^{D+1})\le\deepstrarvol(A),
\]
the arithmetic volume being computed with the chart lattices $\Lambda_\tau(y,a)$.
\formallabel{ReciprocityTheory.ChartRestriction.G.Geo.BirationalPullbackPushforwardFormula}
\end{definition}

\begin{definition}[Resolution in projective coordinates; lines of large degree]\label{def:pf:deepstr:resk}
(a) \emph{Resolution in projective coordinates} $\mathrm{(RPC)}$: for every field $K$ of characteristic $0$ and
every closed subscheme $i\colon X\hookrightarrow\mathbb P^M_K$ which is reduced and has a dense open subset $U$ smooth of
relative dimension $D$, there are $r\in\N$, a closed subscheme $j\colon X'\hookrightarrow\mathbb P^M_K\times\mathbb P^r_K$ and a
morphism $f\colon X'\to X$ with $i\circ f=\mathrm{pr}_1\circ j$ which is a \emph{resolution with coordinates} of
$(X,U)$: $X'$ is smooth of relative dimension $D$, $f$ is an isomorphism over $U$, $f^{-1}(U)$ is dense, and
$X'$ is covered by finitely many affine opens $V$ with coordinates $z_1,\dots,z_D$ such that
$(X'\smallsetminus f^{-1}(U))\cap V$ is the union of the zero loci of $z_1,\dots,z_k$ for some $k\le D$.

(b) \emph{Lines of large degree} $\mathrm{(LLD)}$: for every number field $K$, $i\ge1$, Hermitian lattice
$\overline F$ of rank $i$ over $\OO_K$ and $\eps>0$, there are a finite extension $K'/K$ and $0\ne v\in K'^i$ such
that the normalised Arakelov degree of the line spanned by $v$ in $\overline F\otimes\OO_{K'}$ is at least
$\mu(\overline F)-\frac12\log i-\eps$, $\mu$ the normalised slope (an absolute Siegel lemma in the sense of
\cite{Zhang1995,RoyThunder1996}).
\formal{ReciprocityTheory.ChartRestriction.G.Geo.ResK.ResKStatement}
\formallabel{GroundTheory.Heights₂.Heights.LineOfLargeDegree.Statement}
\end{definition}

\subsubsection*{Derivation of the inputs}

\begin{lemma}[Resolution in projective coordinates]\label{lem:pf:deepstr:resk}
\ref{H:pf:deepstr:res} and \ref{H:pf:deepstr:princ} imply $\mathrm{(RPC)}$.
\formallabel{GlobalTheory.CompanionTheory.Hypotheses₂.ResKSide.resK_of_resolution_principalisation}
\formallabel{PositivityTheory.HtPt.Ah.resK}
\end{lemma}

\begin{proof}
Let $i\colon X\hookrightarrow\mathbb P^M_K$ and $U$ be as in \cref{def:pf:deepstr:resk}(a), and $g\colon X\to\Spec K$
the structure map; $g$ is projective, hence proper. By \ref{H:pf:deepstr:res} (in the projective form, the blow-up
sequence of a projective scheme being projective) there is a sequence of blow-ups $s\colon X_1\to X$ with $X_1$
smooth and projective over $K$, $s$ an isomorphism over the smooth locus of $X$, and the reduced preimage of
the non-smooth locus a divisor with strict normal crossings. As $U$ is smooth, it lies in the smooth locus,
so $s$ is an isomorphism over $U$; the complement of a divisor with strict normal crossings is dense, so
$s^{-1}(U)$ is dense in $X_1$, and $X_1$ is smooth of relative dimension $D$ (it is so on the dense open
$s^{-1}(U)\cong U$). Let $I$ be the ideal sheaf of the reduced preimage of $X\smallsetminus U$; its support
$s^{-1}(X\smallsetminus U)$ contains no irreducible component of $X_1$, since $s^{-1}(U)$ is dense; $X_1$ is
noetherian. By \ref{H:pf:deepstr:princ} applied to $I$ and $F=\varnothing$ there is a sequence of blow-ups
$t\colon X_2\to X_1$ with $X_2$ smooth, $I\OO_{X_2}$ effective Cartier with support of strict normal
crossings, and $t$ an isomorphism over $X_1\smallsetminus\operatorname{supp}I=s^{-1}(U)$. Hence $f=s\circ t$ is an
isomorphism over $U$, $f^{-1}(X\smallsetminus U)=\operatorname{supp}(I\OO_{X_2})$ has strict normal crossings, its
complement $f^{-1}(U)$ is dense, and $X_2$ is projective over $K$ (a sequence of smooth blow-ups of a projective
scheme), smooth of relative dimension $D$. Since $X_2$ is quasi-compact and smooth, $f^{-1}(U)$ is dense and
smooth of relative dimension $D$ and $f^{-1}(X\smallsetminus U)$ has strict normal crossings,
\lref{GlobalTheory.T.coords_statement} shows that $X_2$ is covered by finitely many affine opens $V$ with
coordinates $z_1,\dots,z_D$ such that $f^{-1}(X\smallsetminus U)\cap V$ is the union of the zero loci of
$z_1,\dots,z_k$ for some $k\le D$; this is the condition of adapted coordinates. Finally, $X_2$ being
projective over $K$, it has a closed immersion $e$ into some $\mathbb P^r_K$, and $j=(i\circ f,e)\colon
X_2\to\mathbb P^M_K\times\mathbb P^r_K$ is a closed immersion with $\mathrm{pr}_1\circ j=i\circ f$, because
$\mathrm{pr}_2\circ j=e$ is one and $\mathbb P^M_K$ is separated
(\lref{GlobalTheory.CompanionTheory.Hypotheses₂.exists_closed_in_prod}).
\end{proof}

\begin{lemma}[Full intersection systems]\label{lem:pf:deepstr:ais}
\ref{H:pf:deepstr:res}, \ref{H:pf:deepstr:princ} and \ref{H:pf:deepstr:pl} imply the existence of a full
intersection system. \formallabel{PositivityTheory.HtPt.Ah.ais}
\end{lemma}

\begin{proof}
By \cref{lem:pf:deepstr:resk}, $\mathrm{(RPC)}$ holds. From $\mathrm{(RPC)}$ and \ref{H:pf:deepstr:pl} the
development derives the arithmetic Poincar\'e--Lelong formula in coordinates in its standard form
(independence of the recursion of \cref{def:covolume:hs}(c) from the chosen section),
\lref{GlobalTheory.CompanionTheory.Ar.Tw.Main.poincareLelong_of_inv}; from this formula and
$\mathrm{(RPC)}$, \lref{GlobalTheory.CompanionTheory.Ar.Pub.arithIntersectionSystem_of_poincareLelong_resK}
constructs a full intersection system (by the recursion of Bost, Gillet and Soul\'e
\cite{BostGilletSoule1994}).
\end{proof}

\begin{lemma}[The volume inequality]\label{lem:pf:deepstr:vi}
\ref{H:pf:deepstr:mor} implies $\mathrm{(VI)}$.
\formallabel{GlobalTheory.T.big_line_implies_height_positivity}
\formallabel{GlobalTheory.CompanionTheory.Ar.Small.selfIntersection_le_volNm_of_moriwakiBig}
\end{lemma}

\begin{proof}
Let $P=P_y$, $D\ge1$, $\deepstrII$ and $A=(a,h)$ be as in \cref{def:pf:deepstr:bppf}. A Hermitian line which is
semipositive relative to some $P_0\subseteq P$ is nef on the generic fibre
(\lref{GroundTheory.Heights₂.Heights.HilbertSamuel.NefOnGenericOfSemipos.proof}), so
\ref{H:pf:deepstr:mor} applies and gives $\deepstrII(P;A^{D+1})\le\deepstrarvol'(A)$, where $\deepstrarvol'$ is the
arithmetic volume computed with the groups $\Lambda'_\tau(y,a)\subseteq\Lambda_\tau(y,a)$ of elements $s$ with
$s\prod_ky_{k,c_k}^{-\tau a_k}$ \emph{in} the chart rings (rather than integral over them). By the packing
statement \lref{GlobalTheory.T.packing'} (for a smooth line, the sets of small elements of degree $\tau$ of
$\Lambda_\tau(y,a)$ are finite, and the logarithms of their cardinalities are at most $C(\tau+1)^{D+1}$ for some
$C$), the number of small elements of $\Lambda'_\tau(y,a)$ is at most that of $\Lambda_\tau(y,a)$, and the two
sequences whose $\limsup$ define the volumes are bounded; hence $\deepstrarvol'(A)\le\deepstrarvol(A)$.
\end{proof}

\begin{lemma}[Quotients and the sharp Hodge--Tate statement]\label{lem:pf:deepstr:quot}
\textup{(a)} For every commutative ring $R$, every abelian scheme $A$ over $R$ with a relatively ample line
bundle and every finite flat subgroup scheme $G\subseteq A$, there is a finite flat surjective homomorphism
$q\colon A\to A'$ onto an abelian scheme with a relatively ample line bundle, whose kernel is $G$ (as a functor on
$R$-algebras).

\textup{(b)} \ref{H:pf:deepstr:far} implies the sharp Hodge--Tate statement \ref{G:pf:deepstr:ht}.
\formallabel{PositivityTheory.HtPt.Ah.quot}
\formallabel{GlobalTheory.CompanionTheory.Relay.hodge_tate_of_fargues_exact}
\formallabel{PositivityTheory.Xgn.farguesHodgeTate_of_farguesExact}
\end{lemma}

\begin{proof}
(a) The development proves that quotients of schemes by free actions of finite flat group schemes, with
orbits contained in affine opens, exist and are finite flat
(\lref{GlobalTheory.CompanionTheory.Hypotheses₂.FreeQuotientScheme_holds}), and deduces (a)
(\lref{GlobalTheory.CompanionTheory.Ar.GR.QuotSide.quotientByFiniteFlat_of_freeQuotient}); compare
\cite[\S12]{Mumford1970}. (b) The second half of the sharp Hodge--Tate statement \ref{G:pf:deepstr:ht} is \ref{H:pf:deepstr:far}. For the
first half, \lref{GlobalTheory.CompanionTheory.Ar.HT.hodge_tate_cokernel_sharp} states that
$\kappa\,\omega_{G^D}\subseteq\operatorname{im}(\alpha_G\otimes\OO)$ whenever $v(\kappa)^{p-1}\le v(p)$. For odd $p$ this is the
required statement; for $p=2$ the hypothesis $v(\kappa)\le v(4)=v(2)^2$ implies $v(\kappa)\le v(2)$, since
$v(2)<1$.
\end{proof}

\begin{littheorem}[{\cite[Ch.~I, Prop.~2.7]{FaltingsChai1990}}]\label{lit:pf:deepstr:sah}
Let $S$ be a normal noetherian scheme, $U\subseteq S$ a dense open subscheme and $G,H$ semi-abelian schemes
over $S$. Every homomorphism $G_U\to H_U$ extends uniquely to a homomorphism $G\to H$.
\formallabel{GlobalTheory.CompanionTheory.SAH.semiAbelianHomExtension}
\end{littheorem}

The development proves this statement from the case of a discrete valuation ring.

\begin{theorem}[Lines of large degree]\label{thm:pf:deepstr:lld}
The existence of a full intersection system, $\sfname{HSUse}$, $\sfname{MoriwakiBigStatement}$ and $\mathrm{(RPC)}$
imply $\mathrm{(LLD)}$. \formallabel{GlobalTheory.CompanionTheory.Ar.Giv.lineOfLargeDegree_of_use}
\end{theorem}

\begin{proof}
By \lref{GlobalTheory.CompanionTheory.Ar.Giv.hsLinks_of_use}, $\sfname{MoriwakiBigStatement}$ implies the
successive-minima upper bound: if $A$ is nef on the fibres, semipositive and conjugation-invariant on an
arithmetic variety $P$ of dimension $D+1$, and a Zariski-dense family of arithmetic curves $W\supseteq P$
satisfies $\deepstrII(W;A)\le e'\deg W$, then $\deepstrII(P;A^{D+1})\le(D+1)e'\,\deepstrgdeg(P,A)$; and
$\sfname{HSUse}$ implies the corresponding lower bound (with $e\deg W\le\deepstrII(W;A)$ for all arithmetic
curves $W\supseteq P$) \cite{Zhang1995}. A full
intersection system provides the existence and uniqueness of intersection numbers, their push-forward
formula, a degree formula and coordinate projections of Grassmannian lines
(\lref{GlobalTheory.T.intersection_exists_of_arith_system},
\lref{GlobalTheory.T.generic_degree_of_arith_system},
\lref{GlobalTheory.T.coordinate_projection_of_arith_system}); $\mathrm{(RPC)}$ provides the
integrability and the independence of points (\lref{GlobalTheory.T.points_integrable_of_resolution_statement}).
Together with statements on vertical points, charts, chart energies and dual Grassmannian models proved in
the development, these are the hypotheses of
\lref{GroundTheory.Heights₂.Heights.LineOfLargeDegreeFromZ.ofOpen}, which yields $\mathrm{(LLD)}$.
\end{proof}

\begin{theorem}[Slope bound]\label{thm:pf:deepstr:slopebound}
$\mathrm{(LLD)}$ implies the Lie slope bound (\ref{G:pf:deepstr:slope}).
\formallabel{GlobalTheory.T.slope_bound_of_line_large_degree}
\formallabel{GlobalTheory.CompanionTheory.Ar.Sl.slopeBound_and_lemmaS_of_siegel}
\end{theorem}

\begin{proof}
The formal statement asserts the Lie slope bound twice (the second copy is a synonym used elsewhere in
the development). It is obtained from
\lref{GlobalTheory.T.slope_bound_of_large_degree_line}, which derives the Lie slope bound from
$\mathrm{(LLD)}$ and a bound for the key degree of abelian schemes in terms of their bad reduction, and from
\lref{GlobalTheory.T.key_degree_of_bad_reduction} with
\lref{GlobalTheory.T.key_bad_reduction_holds}, which provide that bound.
\end{proof}

\begin{theorem}[The package $\sfname{Given}$]\label{thm:pf:deepstr:given}
The existence of a full intersection system, $\sfname{HSUse}$, $\sfname{MoriwakiBigStatement}$, $\mathrm{(RPC)}$,
\cref{lit:pf:deepstr:sah}, the sharp Hodge--Tate statement \ref{G:pf:deepstr:ht}, GAGA for sections and $\mathrm{(LLD)}$ imply
$\sfname{Given}$. \formallabel{GlobalTheory.CompanionTheory.Ar.Link.given_of_use}
\formallabel{GlobalTheory.T.final_theorem_of_inputs_left}
\formallabel{GlobalTheory.CompanionTheory.Ar.Top.given_of_inputs}
\end{theorem}

\begin{proof}
The components \ref{G:pf:deepstr:sah}, \ref{G:pf:deepstr:resk}, \ref{G:pf:deepstr:ht}, \ref{G:pf:deepstr:gaga}
are hypotheses, and \ref{G:pf:deepstr:slope} follows from $\mathrm{(LLD)}$ by \cref{thm:pf:deepstr:slopebound}.
The others are proved in the development: \ref{G:pf:deepstr:tor} is
\lref{GlobalTheory.CompanionTheory.Ar.Top.modelClause_holds}; \ref{G:pf:deepstr:key} is
\lref{GlobalTheory.KeyFam.keyFormula_holds}; \ref{G:pf:deepstr:segre} is
\lref{GlobalTheory.CompanionTheory.Segre.numericalSegre_closed}; \ref{G:pf:deepstr:sa} is
\lref{GlobalTheory.T.semi_abelian_statement_holds}; \ref{G:pf:deepstr:anti} is
\lref{GlobalTheory.T.antisymmetric_perturbation_holds}; \ref{G:pf:deepstr:provan} is
\lref{GlobalTheory.T.analytic_provision_holds}; and \ref{G:pf:deepstr:pb}, \ref{G:pf:deepstr:provhs},
\ref{G:pf:deepstr:zhang} follow from the existence of a full intersection system, $\sfname{HSUse}$,
$\sfname{MoriwakiBigStatement}$ and $\mathrm{(RPC)}$ by
\lref{GlobalTheory.CompanionTheory.Ar.Giv.inputs_of_use},
\lref{GlobalTheory.CompanionTheory.Ar.Link.provHS_of_use} and
\lref{GlobalTheory.CompanionTheory.Ar.Link.nineFree_of_use}. (The formal proof first states the
implication with \ref{G:pf:deepstr:tor}, \ref{G:pf:deepstr:key}, \ref{G:pf:deepstr:segre},
\ref{G:pf:deepstr:pb}--\ref{G:pf:deepstr:zhang} as hypotheses and then supplies them.)
\end{proof}

\begin{theorem}\label{thm:pf:deepstr:ofuse}
The existence of a full intersection system, $\sfname{HSUse}$, $\sfname{MoriwakiBigStatement}$, $\mathrm{(RPC)}$,
\cref{lit:pf:deepstr:sah}, the sharp Hodge--Tate statement \ref{G:pf:deepstr:ht}, GAGA for sections, the existence of quotients by finite
flat subgroups (\cref{lem:pf:deepstr:quot}(a)), \ref{H:pf:deepstr:zar} and \ref{H:pf:deepstr:nef} imply
$\sfname{DeepStR}$. \formallabel{GlobalTheory.CompanionTheory.Ar.Plug.deepStR_of_use}
\end{theorem}

\begin{proof}
By \cref{thm:pf:deepstr:lld}, $\mathrm{(LLD)}$ holds, hence $\sfname{Given}$ by \cref{thm:pf:deepstr:given}.
The good tower volume theorem \lref{GlobalTheory.CompanionTheory.Ar.Link.good_tower_volume_theorem_bis}
(\cref{thm:tower}), whose hypotheses are the existence of a full intersection system, $\sfname{HSUse}$,
$\sfname{MoriwakiBigStatement}$, $\mathrm{(RPC)}$, \cref{lit:pf:deepstr:sah}, the sharp Hodge--Tate statement \ref{G:pf:deepstr:ht},
GAGA for sections, $\mathrm{(LLD)}$, the quotient statement and \ref{H:pf:deepstr:zar}, gives
$\sfname{GoodTowerExistsX}$ and $\sfname{VolumeTheoremX}$. The statement $\mathrm{(VI)}$ is the fifth
component of the Hilbert--Samuel provisions \ref{G:pf:deepstr:provhs} (it also follows directly from
\ref{H:pf:deepstr:mor} by \cref{lem:pf:deepstr:vi}, which is how the development proves that component). Now \cref{thm:pf:deepstr:ofvolume} below,
applied with \ref{H:pf:deepstr:nef}, $\sfname{Given}$, the full intersection system, $\mathrm{(VI)}$ and the
two tower propositions, gives $\sfname{DeepStR}$.
\end{proof}

\begin{proof}[Proof of \cref{pf:deepstr:main}]
By \cref{lem:pf:deepstr:ais,lem:pf:deepstr:resk}, \ref{H:pf:deepstr:res}--\ref{H:pf:deepstr:pl} give a full
intersection system and $\mathrm{(RPC)}$. \Cref{lit:pf:deepstr:sah} holds; GAGA for sections is proved in
the development (\lref{GlobalTheory.CompanionTheory.Gaga.gagaSections}); the quotient statement is
\cref{lem:pf:deepstr:quot}(a); and the sharp Hodge--Tate statement \ref{G:pf:deepstr:ht} follows from \ref{H:pf:deepstr:far} by
\cref{lem:pf:deepstr:quot}(b). With \ref{H:pf:deepstr:hs}, \ref{H:pf:deepstr:mor}, \ref{H:pf:deepstr:zar}
and \ref{H:pf:deepstr:nef}, \cref{thm:pf:deepstr:ofuse} gives $\sfname{DeepStR}$.
\end{proof}

% ---------------------------------------------------------------------------
\subsection{From the volume theorem to \texorpdfstring{$\sfname{DeepStR}$}{DeepStR}}\label{sec:pf:deepstr:volume}

\subsubsection*{The uniform mean-slope bound}

\begin{theorem}[Mean slopes on unitary moduli spaces]\label{thm:pf:deepstr:meanall}
Assume \ref{H:pf:deepstr:nef}, $\sfname{Given}$, the existence of a full intersection system, $\mathrm{(VI)}$,
$\sfname{GoodTowerExistsX}$ and $\sfname{VolumeTheoremX}$. There is $C\ge0$ with the following property. Let
$E$ be an abelian CM field with host $H$, $D_w$ a datum of signature $(p,q)$, $l$ a level, $X^o$ a moduli
space, $c$ a coefficient datum, $r\in\N$ and $U$ uniformisations with the composition laws $\mathrm{(CL)}$ for some
$h,e$, such that $\mathrm{(UP)}$ and $\mathrm{(GI)}$ hold, $l$ has good reduction at every prime,
$l_A=1$, $q\ge1$, $3q+1\le p+q$, $|\Psi_H|\ge2$ and
\begin{equation}\label{eq:pf:deepstr:badhyp}
\#\{p\text{ prime}:\ p\mid2|d_E|\ \text{or}\ p\mid M'\}\le4|\Psi_H|(p+q)^2.
\end{equation}
Let $x\in\R$ with $|\tau(\xi_0)|=x$ for the base points $\xi_0$ of all $U_\sigma$ with $\sigma$ good and all
$\tau\colon E\to\C$. Then $\deepstrMeanAll(\log\deepstrrd_E+|L_H|,C)$ holds for $(X^o,c,U,x)$.
\formallabel{GlobalTheory.T.abelian_cm_comp_laws_bound}
\formallabel{GlobalTheory.CompanionTheory.Ar.Plug.top_meanAll}
\end{theorem}

The proof uses three results proved below: a uniform version of the tower propositions
(\cref{prop:pf:deepstr:vstower}), the comparison theorem (\cref{thm:pf:deepstr:comparison}) and the bound
for the cost of clearing denominators (\cref{lem:pf:deepstr:costbound}).

\begin{proposition}\label{prop:pf:deepstr:vstower}
Assume $\sfname{Given}$, $\sfname{GoodTowerExistsX}$ and $\sfname{VolumeTheoremX}$. There is $C_T\ge0$ such
that for all $E,H,D_w,l,X^o,c$ with $E$ abelian, every full intersection system $\deepstrII$, and under the
hypotheses $\mathrm{(GI)}$, good reduction of $l$ at every prime, $l_A=1$, $q\ge1$, $3q+1\le p+q$ and
$\#\{p\text{ prime}:p\mid2d_EM'\}\le4|\Psi_H|(p+q)^2$, there are a model $Y$ of $X^o$ over $\OO_{F_c}$, a root
tower $T$ of type $(\Psi_H,p+q)$, a good origin $o$ of $T$ on $Y$ and a bijection
$\kappa\colon\deepstrComp(X^o,c)\to\deepstrComp(T)$ such that $o$ satisfies $\mathrm{(IC)}$, $\kappa$ is compatible
with $o$, $\mathrm{bad}(T)$ is the set of
prime divisors of $2d_EM'$, and $T$ satisfies proportionality, $\mathrm{(LC)}$, equality of volumes and
$\mathrm{(MS)}$ with constants $(C_T,\deepstrrd_E,L_H)$.
\formallabel{GlobalTheory.T.full_intersection_good_level}
\formallabel{GlobalTheory.CompanionTheory.Ar.Plug.vs_tower}
\end{proposition}

\begin{proof}
Let $C_T$ be the constant of $\sfname{VolumeTheoremX}$. Given the data, note $p+q\ge1$ and $p\ge1$ (as
$3q+1\le p+q$). Choose for every level $l'$ a moduli space $X^o_{l'}$, with $X^o_l=X^o$; they exist by
\lref{GlobalTheory.Moduli.moduli_scheme_exists}. Since $l$ has good reduction everywhere and $l_A=1$,
the toroidal package \ref{G:pf:deepstr:tor} gives a global model $Y$ of $X^o$ over $\OO_{F_c}$
(\lref{GlobalTheory.CompanionTheory.Ar.Rec.exists_isGlobal}, which states that a global model exists for
every level with good reduction everywhere and trivial $A$-part). By $\mathrm{(GI)}$ the components of $Y$ are split
(\lref{GlobalTheory.T.split_components_of_alg_closed}, which deduces split components of every model
from (GI)), so $\sfname{GoodTowerExistsX}$ gives a root tower
$T$ and a good origin $o$ with $\mathrm{(IC)}$, and \lref{GlobalTheory.T.tower_components_equiv_of_given}, which
states that under $\sfname{Given}$ every tower origin admits a compatible bijection of components, gives
$\kappa$.

By definition of a tower origin, $r\in\mathrm{bad}(T)$ if and only if $r$ is a prime dividing $2d_EM'$. If
$2d_EM'$ were $0$, every prime would lie in the finite set $\mathrm{bad}(T)$, which is absurd; hence
$2d_EM'\ne0$ and $\mathrm{bad}(T)$ is its set of prime divisors. In particular
$\#\mathrm{bad}(T)\le4|\Psi_H|(p+q)^2$, and $\sfname{VolumeTheoremX}$ (applicable as $E$ is abelian, $Y$ is
global, $\mathrm{(IC)}$ holds, $\deepstrII$ is full, $q\ge1$ and $3q+1\le p+q$) gives proportionality,
equality of volumes, $\mathrm{(LC)}$ and $\mathrm{(MS)}$ with constants $(C_T,\deepstrrd_E,L_H)$.
\end{proof}

\begin{lemma}[Integer weights]\label{lem:pf:deepstr:intweights}
Let $u$ be an integer weight for a datum of signature $(p,q)$, $N=p+q$, and let $\lambda=u$ viewed as a real
weight. Then $\|u\|_1=\|\lambda\|_1$, $\deepstrbnd(u,g)=\deepstrbnd_0(\lambda,g)$ for all $g$, and if $1\le q<p$
and $\deepstrWtGap(u,g)$ holds, then $\mathrm{WH}(p,\lambda,g)$ holds.
\formallabel{GlobalTheory.T.tower_weight_of_zero}
\formallabel{GlobalTheory.CompanionTheory.Ar.Plug.integer_weights_vocab}
\end{lemma}

\begin{proof}
The first two identities hold by definition ($\|\lambda\|$ is the maximum of the $|u_{\tau,i}|$). Let
$1\le q<p$ and $\deepstrWtGap(u,g)$. Then $g>0$, and $\lambda$ is regular dominant because each $u_\tau$ is
strictly decreasing. As $j=p>q=N-j$, the small block of $j=p$ is $\{i\ge p\}$, so the last condition of
$\mathrm{WH}$ is the gap condition of $\deepstrWtGap$. Finally the gap of $\lambda$ at $\tau$ is
$u_{\tau,p-1}-u_{\tau,p}=|u_{\tau,p}-u_{\tau,p-1}|\ge g$, by $\deepstrWtGap$ with $a=p$ (which exists as
$q\ge1$) and $b=p-1$.
\end{proof}

The formal statement of \cref{lem:pf:deepstr:intweights} also contains four identities relating $u$ to the
coefficient vectors of root models, to the Vandermonde product and to the automorphic weight $G_u$, which are
not used in this section.

\begin{proof}[Proof of \cref{thm:pf:deepstr:meanall}]
Fix a full intersection system $\deepstrII$, let $C_T$ be as in \cref{prop:pf:deepstr:vstower} and put
$C=C_T+120$. Let the data be as in the statement, and write $\Psi=\Psi_H$, $N=p+q$, $\deepstrrd=\deepstrrd_E$,
$L=L_H$. Then $q<p$ and $N\ge4$ (from $q\ge1$ and $3q+1\le N$). The set of prime divisors of $2d_EM'$ is
contained in the set in \eqref{eq:pf:deepstr:badhyp}, so \cref{prop:pf:deepstr:vstower} applies and gives
$Y,T,o,\kappa$; we have $j=p$ for $T$, and $\min(j,N-j)=q$.

Let $(u,g)$ satisfy $\deepstrWtGap(u,g)$ and put $\lambda=u$ (real), $B=\deepstrbnd(u,g)$ and $B_0=\lfloor B\rfloor$.
By \cref{lem:pf:deepstr:intweights}, $\mathrm{WH}(p,\lambda,g)$ holds and $B=\deepstrbnd_0(\lambda,g)$. As
$|\Psi|\ge2$ and $N\ge4$, $B\ge4|\Psi|N^2\ge128$, so $B_0\ge128$, $B<B_0+1$, and $4|\Psi|N^2\le B_0$ (since
$4|\Psi|N^2$ is an integer $\le B$). Hence $\#\mathrm{bad}(T)\le4|\Psi|N^2\le B_0$, and
$\mathrm{(MS)}$ at $B_0$ gives
\begin{equation}\label{eq:pf:deepstr:hM}
S:=\frac1{\#\deepstrComp(T)}\sum_{c}\deepstrslope_T(B_0,c,\lambda)\ge-C_T\bigl(\Lambda(B_0,\deepstrrd,L)\|u\|_1+q\,\deepstrspread(u)\bigr).
\end{equation}
For $Z\in\deepstrComp(X^o,c)$ put
$\mathrm{sl}(Z)=\deepstrslope_T(B_0,\kappa(Z),\lambda)-2\deepstrcost(l,B)\|u\|_1$, with $\deepstrcost(l,B)$ as in
\cref{def:pf:deepstr:cost}. By \cref{thm:pf:deepstr:comparison} (applicable: $E$ is Galois over $\Q$, each
$u_\tau$ is strictly decreasing, $\mathrm{WH}(p,\lambda,g)$ holds and $\deepstrbnd_0(\lambda,g)\le B$),
$\mathrm{sl}(Z)\le\deepstrvsl(Z,u)$ for the valuation family $\mathcal V_B$. As $\kappa$ is a bijection,
\[
\frac1{\#\deepstrComp(X^o,c)}\sum_Z\mathrm{sl}(Z)=S-2\deepstrcost(l,B)\|u\|_1.
\]
By \cref{lem:pf:deepstr:costbound} (applicable since $B\ge128$ and the set in \eqref{eq:pf:deepstr:badhyp} has at
most $4|\Psi|N^2\le B_0$ elements), $0\le\deepstrcost(l,B)\le60(\log B+\log\deepstrrd)$. Moreover
$\Lambda(B_0,\deepstrrd,L)=\log B_0+\log\deepstrrd+|L|\le\log B+\log\deepstrrd+|L|=:\Lambda'$, and
$\Lambda'\ge0$ because $B\ge1$ and $\deepstrrd\ge1$. Hence, with $\|u\|_1\ge0$ and $q\,\deepstrspread(u)\ge0$,
\begin{align*}
S-2\deepstrcost(l,B)\|u\|_1&\ge-C_T\bigl(\Lambda'\|u\|_1+q\,\deepstrspread(u)\bigr)-120(\log B+\log\deepstrrd)\|u\|_1\\
&\ge-(C_T+120)\bigl(\Lambda'\|u\|_1+q\,\deepstrspread(u)\bigr)=-C\,\deepstrerr_{\log\deepstrrd+|L|}(u,g),
\end{align*}
the last equality being the definition of $\deepstrerr$ with $\log\deepstrbnd(u,g)=\log B$. This is
$\deepstrMeanGe(u,B,C\,\deepstrerr_{\log\deepstrrd+|L|}(u,g))$.
\end{proof}

\subsubsection*{Deep covers}

\begin{lemma}[Parameters of deep covers]\label{lem:pf:deepstr:params}
Assume $\sfname{Given}$. Let $\mathcal D$ be a global datum with ranks $k,m$, $n\ge1$ and $f=2n$ as in
\cref{sec:families}, $r=k+m-1$, $T$ a core datum of $\mathcal D$ satisfying $\mathsf{LawsJ}$, $1\le j<m$, and $Y$
a deep cover of $T$ at level $j$ which is good for some parameters. Then the uniformisations of $Y$ satisfy
the composition laws $\mathrm{(CL)}$ for suitable generic frames, the level of $Y$ has good reduction at
every prime, the datum $D_w(j)$ of $T$ at level $j$ has signature $(p,q)$ with $q=m-j$ and $p+q=r$,
$|\Psi_H|=f$, $\deepstrrd_E$ is the root discriminant of the host field of $\mathcal D$, and
$L_H=\Re\frac{L'}{L}(1,\chi_d)$ is the value used in the definition of $\ell(\mathcal D)$
(\cref{def:covolume:deepstr}).
\formallabel{GlobalTheory.T.good_datum_parameters}
\formallabel{GlobalTheory.CompanionTheory.Ar.Plug.vs_laws}
\end{lemma}

\begin{proof}
The laws $\mathsf{LawsJ}$ fix the signature of $D_w(j)$ to be $(k+j-1,m-j)$, whence $q=m-j$ and
$p+q=k+m-1=r$. They also make the modulus of the level $l$ of $T$ the product $g_0g_1$ of the two type primes
of the host; \lref{Formalism.L2.Pj.goodAt_of_two_primes} states that a level whose modulus is a product
of two distinct primes has good reduction at every prime. Goodness of $Y$ makes the modulus of the level of $Y$ a multiple of that of $l$, and
\lref{GlobalTheory.LawsDeep.goodAt_of_mul} states that good reduction at every prime passes to such a
level. The composition laws follow from the existence statements for moduli spaces and uniformisations in
\ref{G:pf:deepstr:tor} and the conditions imposed by goodness on the uniformisations of $Y$
(\lref{GlobalTheory.CompanionTheory.Ar.Plug.VsLaws.compLaws_of}, which derives (CL) from these data).
The identity $|\Psi_H|=f$ is a property of the CM type of the host,
\lref{GlobalTheory.CompanionTheory.Ar.Plug.VsLaws.hostLval_eq} identifies $L_H$ with the logarithmic
derivative used in $\ell(\mathcal D)$, and the identity for $\deepstrrd_E$ holds by definition.
\end{proof}

\begin{theorem}\label{thm:pf:deepstr:ofvolume}
Assume \ref{H:pf:deepstr:nef}, $\sfname{Given}$, the existence of a full intersection system, $\mathrm{(VI)}$,
$\sfname{GoodTowerExistsX}$ and $\sfname{VolumeTheoremX}$. Then $\sfname{DeepStR}$ holds.
\formallabel{GlobalTheory.CompanionTheory.Ar.Plug.deepStR_of_volume}
\end{theorem}

\begin{proof}
Let $C$ be the constant of \cref{thm:pf:deepstr:meanall}. Let $\mathcal D,T,j,Y$ be as in $\sfname{DeepStR}$:
$T$ satisfies $\mathsf{LawsJ}$, $m^2\le k$, $1\le j<m$, $Y$ is good for some parameters, and the level modulus
$M'$ of $Y$ has at most $15$ prime factors. The host field $E$ of $\mathcal D$ is abelian over $\Q$
(\lref{GlobalTheory.CompanionTheory.Ar.Rec.field_of_datum_abelian}: it is Galois over $\Q$ with commutative
Galois group). By \cref{lem:pf:deepstr:params}, the datum of $Y$ has $q=m-j$ and $p+q=r=k+m-1$, and
$|\Psi_H|=2n$. We check the hypotheses of \cref{thm:pf:deepstr:meanall}.
\begin{itemize}
\item $q=m-j\ge1$ since $j<m$.
\item $3q+1\le p+q$: as $m\ge2$, $k\ge m^2\ge2m$, so $3q+1=3m-3j+1\le3m-2\le k+m-1$.
\item $|\Psi_H|=2n\ge2$ since $n\ge1$.
\item \eqref{eq:pf:deepstr:badhyp}: let $\mathcal S$ be the set of primes dividing $2|d_E|$ or $M'$.
\Lref{Formalism.BasicNotions.Host.card_bad_gt_le} states that for every finite set $S'$ and every
$B\ge n$ at most $\#S'+3$ primes $p>B$ divide $2|d_E|$ or lie in $S'$ (the prime divisors of $2|d_E|$ lie in the
set of small primes of the host, at most one of which exceeds $n$, together with its distinguished primes
$q,d$; \cref{sec:families}). With $S'$ the set of prime factors of $M'$ and $B=n$, at most $15+3$ elements of
$\mathcal S$ exceed $n$, and at most $n+1$ elements are $\le n$. So
$\#\mathcal S\le n+19\le128n=4\cdot2n\cdot16\le4|\Psi_H|(p+q)^2$, using $n\ge1$ and $p+q=k+m-1\ge4$.
\end{itemize}
The remaining hypotheses are supplied by the deep cover: (CL) and good reduction everywhere by
\cref{lem:pf:deepstr:params}; (UP) by \lref{GlobalTheory.CompanionTheory.Ar.Ref.unifAll_holds}, which
states that (UP) holds for every moduli space and coefficient datum; (GI), $l_A=1$ and the normalisation
$|\tau(\xi_0)|=\xi_R(\mathcal D)$ of the base points by the conditions of goodness. Hence
$\deepstrMeanAll(\log\deepstrrd_E+|L_H|,C)$ holds on $Y$ with $x=\xi_R(\mathcal D)$, and
$\log\deepstrrd_E+|L_H|=\ell(\mathcal D)$ by \cref{lem:pf:deepstr:params}.
\end{proof}

% ---------------------------------------------------------------------------
\subsection{The comparison theorem}\label{sec:pf:deepstr:comparison}

\subsubsection*{The cost of clearing denominators}

\begin{definition}[Cost]\label{def:pf:deepstr:cost}
Let $E$ be a number field, $l$ a level with modulus $M'$ and $B>0$. For a prime $p$ put
\[
\deepstrex(p)=\frac1{p-1}+\frac{40\,v_p(|d_E|)}{[E:\Q]}+\begin{cases}1&p=2,\\0&p\ne2,\end{cases}\qquad
\deepstrcost(l,B)=\sum_{p\in\deepstrPrimes(l,B)}\deepstrex(p)\log p,
\]
with $\deepstrPrimes(l,B)$ as in \eqref{eq:pf:deepstr:htprimes}. For a finite set $S$ of primes and an
integer weight $u$ put $\deepstrclear(S,u)=\prod_{p\in S}p^{\lceil2\deepstrex(p)\|u\|_1\rceil}\in\Z_{\ge1}$.
Finally, the \emph{valuation comparison} $\mathrm{(VC)}$ for $(X^o,c)$ is the condition: for every prime $p$,
every height-one prime $\lambda\ni p$ of $\OO_{F_c}$ and every $w\in\mathcal V^{\mathrm{ht}}(\lambda)$ there is
$v\in\mathcal V^{\mathrm{rec}}(\lambda)$ such that $v(f)\le w(f)+v_\lambda(p)\,\deepstrex(p)\,\|u\|_1$ for every
integer weight $u$ and every section $f$ of weight $G_u$.
\formallabel{GlobalTheory.CompanionTheory.Ar.Plug.exTwo}
\formallabel{GlobalTheory.CompanionTheory.Ar.Plug.costTwo}
\formallabel{GlobalTheory.CompanionTheory.Ar.Plug.SlopeVol.clearN}
\formallabel{GlobalTheory.CompanionTheory.Ar.Plug.HTLow}
\end{definition}

The first two terms of $\deepstrex(p)$ form the exponent of the primes in the covolume bound
(\lref{GlobalTheory.CovHt.htEx} defines it by this formula).

\begin{lemma}[A Mertens-type bound]\label{lem:pf:deepstr:mertens}
For every integer $N\ge1$, $\displaystyle\sum_{p\le N}\frac{\log p}{p-1}\le\log N+2.2$.
\formallabel{ReciprocityTheory.ChartRestriction.G.H1.primeLogSum_nat}
\end{lemma}

\begin{proof}
For $N=1$ the sum is empty. Let $N\ge2$. Write $\frac{\log p}{p-1}=\frac{\log p}p+\frac{\log p}{p(p-1)}$.
For $p\le N$, Legendre's formula gives $v_p(N!)\ge\lfloor N/p\rfloor\ge N/p-1$, so
$N\frac{\log p}p\le v_p(N!)\log p+\log p$; summing over $p\le N$,
\[
N\sum_{p\le N}\frac{\log p}p\le\log N!+\sum_{p\le N}\log p\le(N+1)\log N-N+1+N\log4,
\]
where $\log N!\le(N+1)\log N-N+1$ (as $\log k\le\int_k^{k+1}\log x\,dx$, $\sum_{k=2}^{N-1}\log k\le\int_1^N\log x\,dx=N\log N-N+1$) and $\prod_{p\le N}p\le4^N$ (the classical bound,
proved by induction from $\prod_{m+1<p\le2m+1}p\mid\binom{2m+1}m\le4^m$). Since $\log N\le\log2+N/2-1$
(from $\log x\le x-1$ at $x=N/2$),
\[
\sum_{p\le N}\frac{\log p}p\le\log N+\frac{\log N+1}N-1+\log4\le\log N-\tfrac12+\log4+\frac{\log2}N
\le\log N+1.2329.
\]
For the second part, bound $\log p\le k\log2$ for $p\le2^k$; the primes $p\le30$ contribute at most
$\log2\,(\frac12+\frac26+\frac3{20}+\frac3{42}+\frac4{110}+\frac4{156}+\frac5{272}+\frac5{342}+\frac5{506}+\frac5{812})
\le0.8081$, and for the integers $n\ge31$, by summation by parts with $\frac1{n(n-1)}=\frac1{n-1}-\frac1n$ and
$\log(n+1)-\log n\le\frac1n$,
\[
\sum_{n=31}^N\frac{\log n}{n(n-1)}\le\frac{\log31}{30}+\sum_{n\ge31}\frac1{n^2}\le\frac{5\log2+1}{30}\le0.1489.
\]
The total is at most $\log N+1.2329+0.8081+0.1489\le\log N+2.2$.
\end{proof}

\begin{lemma}[Cost bound]\label{lem:pf:deepstr:costbound}
Let $E$ be a CM field with host $H$, $l$ a level and $B\ge128$ such that at most $\lfloor B\rfloor$ primes divide
$2|d_E|$ or $M'$. Then $0\le\deepstrcost(l,B)\le60(\log B+\log\deepstrrd_E)$.
\formallabel{GlobalTheory.T.cost_two_log_bound}
\formallabel{GlobalTheory.CompanionTheory.Ar.Plug.costTwo_bound}
\end{lemma}

\begin{proof}
Let $N=\lfloor B\rfloor\ge128$, $\mathcal P=\{p\le N\}$ and $\mathcal Q$ the set of primes dividing $2|d_E|$ or
$M'$, so that $\deepstrPrimes(l,B)=\mathcal P\cup\mathcal Q$ (an integer $p$ satisfies $p\le B$ if and only if
$p\le N$). Put $n=[E:\Q]$. Then
\[
\deepstrcost(l,B)=\sum_{p\in\mathcal P\cup\mathcal Q}\frac{\log p}{p-1}
+\frac{40}n\sum_{p\in\mathcal P\cup\mathcal Q}v_p(|d_E|)\log p+\log2,
\]
the last term because $2\in\mathcal Q$. All terms are non-negative, whence $\deepstrcost(l,B)\ge0$. Every prime
divisor of $d_E$ lies in $\mathcal Q$, so the middle sum is $\frac{40}n\log|d_E|=40\log\deepstrrd_E$, and
$\log\deepstrrd_E\ge0$. By \cref{lem:pf:deepstr:mertens} the primes of $\mathcal P$ contribute at most $\log N+2.2$ to
the first sum. Each prime $p\in\mathcal Q\smallsetminus\mathcal P$ satisfies $p\ge N+1$, so
$\frac{\log p}{p-1}\le\frac{\log(N+1)}N$ (the function $x\mapsto\log x/(x-1)$ decreases on $x>1$), and there are at most
$N$ of them by hypothesis; they contribute at most $\log(N+1)\le\log N+1$. Hence
\[
\deepstrcost(l,B)\le2\log N+3.2+\log2+40\log\deepstrrd_E\le2\log B+3.9+40\log\deepstrrd_E
\le60(\log B+\log\deepstrrd_E),
\]
since $\log B\ge7\log2>4.8$.
\end{proof}

\begin{lemma}[Clearing denominators]\label{lem:pf:deepstr:clear}
\textup{(a)} For a finite set $S$ of primes and an integer weight $u$,
$\deepstrclear(S,u)\le\exp\bigl(2\|u\|_1\sum_{p\in S}\deepstrex(p)\log p+\sum_{p\in S}\log p\bigr)$.

\textup{(b)} Assume $\mathrm{(VC)}$, let $B>0$ and let $S=\deepstrPrimes(l,B)$ (a finite set). Let $f$ be a cusp
form of weight $G_u$ on a component $Z$ such that $v(f)\ge-v_\lambda(p)\,\deepstrex(p)\,\|u\|_1$ for every prime
$p\in S$, every $\lambda\ni p$ and every $v\in\mathcal V^{\mathrm{rec}}(\lambda)$, and $v(f)\ge0$ for every prime
$p\notin S$, every $\lambda\ni p$ and every $v\in\mathcal V^{\mathrm{rec}}(\lambda)$. Then $w(\deepstrclear(S,u)f)\ge0$
for every $\lambda$ and every $w\in\mathcal V_B(\lambda)$.
\formallabel{GlobalTheory.CompanionTheory.Ar.Plug.SlopeVol.clearN_le}
\formallabel{GlobalTheory.CompanionTheory.Ar.Plug.SlopeVol.clear_integral}
\end{lemma}

\begin{proof}
(a) For $p\in S$, $p^{\lceil a\rceil}\le p^{a+1}=\exp((a+1)\log p)$ with $a=2\deepstrex(p)\|u\|_1\ge0$; take the
product. (b) Put $N=\deepstrclear(S,u)$; then $w(Nf)=v_\lambda(N)+w(f)$. Let $\lambda$ lie over the prime $p$.
If $p\in S$, then $\mathcal V_B(\lambda)=\mathcal V^{\mathrm{ht}}(\lambda)$, and $\mathrm{(VC)}$ gives
$v\in\mathcal V^{\mathrm{rec}}(\lambda)$ with $v(f)\le w(f)+v_\lambda(p)\deepstrex(p)\|u\|_1$; with the hypothesis,
$w(f)\ge-2v_\lambda(p)\deepstrex(p)\|u\|_1$, while $v_\lambda(N)\ge\lceil2\deepstrex(p)\|u\|_1\rceil v_\lambda(p)\ge
2v_\lambda(p)\deepstrex(p)\|u\|_1$, the other factors of $N$ having non-negative valuation. If $p\notin S$, then
$\mathcal V_B(\lambda)=\mathcal V^{\mathrm{rec}}(\lambda)$, $w(f)\ge0$ by hypothesis and $v_\lambda(N)\ge0$. In both cases
$w(Nf)\ge0$ (trivially if $w(f)=+\infty$).
\end{proof}

\subsubsection*{Counting lemmas}

\begin{lemma}\label{lem:pf:deepstr:counting}
\textup{(a)} Let $\mathcal A,\mathcal B$ be sets, $d\ge1$, and $\mathsf P(a,i,b)$ a relation ($a\in\mathcal A$,
$1\le i\le d$, $b\in\mathcal B$) such that every $a$ and $i$ admit some $b\in\mathcal B$ with $\mathsf P(a,i,b)$, and
$a=a'$ whenever for every $i$ some $b$ satisfies $\mathsf P(a,i,b)$ and $\mathsf P(a',i,b)$. Then
$\#\mathcal A\le(\#\mathcal B)^d$ (cardinalities in $\N\cup\{\infty\}$).

\textup{(b)} Let $\Omega$ be a field, $\Lambda\subseteq\Omega$ a subgroup, $\mathcal S\subseteq\Omega$, $\theta\in\Omega^\times$, and $\mathsf Q$ a
property of elements of $\Lambda$ such that $\sigma\theta\in\mathcal S$ whenever $\mathsf Q(\sigma)$. Then the number of
$\sigma\in\Lambda$ with $\mathsf Q(\sigma)$, taken to be $0$ if infinite, is at most $\#\mathcal S$.

\textup{(c)} Let $D\in\N$, $n,d\ge1$, $k_0\in8\N$, and
$m,m'\colon\N\to\N\cup\{\infty\}$ with $m(\tau)\le m'(\tau+k_0)^d$ for all large $\tau$. Then
$\deepstrcvol(D,nd,8;m)\le\deepstrcvol(D,n,8;m')$.
\formallabel{GlobalTheory.CompanionTheory.Ar.Plug.SlopeVol.encard_le_pow_of_rel}
\formallabel{GlobalTheory.CompanionTheory.Ar.Plug.card_le_encard_of_mul}
\formallabel{GlobalTheory.T.count_volume_power_comparison}
\formallabel{GlobalTheory.CompanionTheory.Ar.Plug.countVol_le_of_le_shift_pow}
\end{lemma}

\begin{proof}
(a) Choose $b(a,i)\in\mathcal B$ with $\mathsf P(a,i,b(a,i))$; the map $a\mapsto(b(a,i))_i\in\mathcal B^d$ is injective
by the second hypothesis. (b) Multiplication by $\theta$ is injective. (c) Write $k_0=8j$. For large $k$,
$m(8k)\le m'(8(k+j))^d$, hence $\deepstrlogc(m(8k))\le d\,\deepstrlogc(m'(8(k+j)))$ (both sides are
monotone in their arguments and $\deepstrlogc(x^d)=d\,\deepstrlogc(x)$), and
\[
\frac{D!\,\deepstrlogc(m(8k))}{nd\,(8k)^{D+1}}\le b_{k+j}\Bigl(1+\frac jk\Bigr)^{D+1},\qquad
b_k=\frac{D!\,\deepstrlogc(m'(8k))}{n\,(8k)^{D+1}}.
\]
As $(1+j/k)^{D+1}\to1$, the $\limsup$ of the right-hand side is at most $\limsup_kb_{k+j}=\limsup_kb_k$
(the product of a sequence in $[0,\infty]$ with a sequence tending to $1$; the case $\limsup b=\infty$ being
trivial).
\end{proof}

\subsubsection*{Geometric inputs}

\begin{proposition}[Inputs on root models and origins]\label{prop:pf:deepstr:geominputs}
The following statements are proved in the development; each item cites its formal proof. Let $E$ be a CM field Galois over $\Q$ with host $H$, $D_w,l,X^o,c$ as before, $U$ uniformisations with
the composition laws $\mathrm{(CL)}$, $x$ with $|\tau(\xi_0)|=x$ as in \cref{thm:pf:deepstr:meanall}, $Y$ a model, $T$ a
root tower with a good origin $o$, $\kappa$ a bijection compatible with $o$, and $\deepstrII$ a full intersection
system such that $T$ is proportional and satisfies $\mathrm{(LC)}$. Assume $\sfname{Given}$. For
$B\in\N$ and $b\in\deepstrComp(T)$ write $\mathrm{Or}=\mathrm{Or}(B,b)$, $M=M(B,b)$, $d=\deg\mathrm{Fl}$.
\begin{enumerate}[label=(\alph*)]
\item\label{gi:topdeg} If \ref{H:pf:deepstr:nef} holds, then for every component $Z$ and every integer weight $u$
with strictly decreasing $u_\tau$: $\mathrm{gr}(Z,G_u)=\deepstrtop(Z,u)=\deepstrvol_T(\kappa(Z))\,\deepstrflagdeg(u)>0$
(\lref{GlobalTheory.T.uniform_norm_degree_bound}).
\item\label{gi:frame} For every $B$ and $Z$, with $b=\kappa(Z)$, there are a generic frame $\mathrm{Fr}$ of
$\mathrm{Or}$ and a reading $\mathrm{Rd}$ of it on $Z$ (\cref{def:pf:deepstr:origins}(e)) which is
\emph{geometric}: the image under $\varphi$ of the generic frame of $Z$ differs from $\mathrm{Fr}$ by the units of
the reading (\lref{GlobalTheory.T.tower_origin_component_bound}).
\item\label{gi:basis} For $(\mathrm{Fr},\mathrm{Rd})$ as in \ref{gi:frame}, every integer weight $u$ and $g,B$ with
$\mathrm{WH}(j,u,g)$ and $\deepstrbnd_0(u,g)<B+1$, there are $k_0\in8\N$, $R_0>0$ and
$f_1,\dots,f_d\in\deepstrmsmall(\mathrm{Or},k_0u,R_0)$, linearly independent over $\varphi\Omega(Z)$, such that for every
integer weight $v$ and $s\in\Omega_M$: if $\deepstrtr_v(sf_i)=0$ for all $i$, then $s=0$
(\lref{GlobalTheory.T.generic_frame_component_bound}).
\item\label{gi:htlow} $\mathrm{(VC)}$ holds (\lref{GlobalTheory.T.low_height_of_given}).
\item\label{gi:mul} $\deepstrmsmall(\mathrm{Or},u,R)\cdot\deepstrmsmall(\mathrm{Or},v,R')\subseteq\deepstrmsmall(\mathrm{Or},u+v,RR')$ for
$R,R'\ge0$ (\lref{GlobalTheory.T.model_small_sections_mul}).
\item\label{gi:descent} \emph{Descent.} If $\mathrm{(UP)}$ holds and $(\mathrm{Fr},\mathrm{Rd})$ is as in
\ref{gi:frame}, then for every integer weight $u$ with strictly decreasing $u_\tau$, every $R\ge0$ and every
$s\in\deepstrmsmall(\mathrm{Or},u,R)$ there is $f\in\deepstrcusp(Z,G_u)$ with $\deepstrev(f)=\deepstrtr_u(s)$,
$\|f\|_{G_u,\sigma,z}\le dR$ for all good $\sigma$ and frames $z$ on $Z$, and
$v(f)\ge-v_\lambda(p)\,[p\in S(B)]\,\deepstrex(p)\,\|u\|_1$ for all primes $p$, all $\lambda\ni p$ and all
$v\in\mathcal V^{\mathrm{rec}}(\lambda)$ (\lref{GlobalTheory.T.uniform_norm_all_tower_bound}).
\item\label{gi:approx} \emph{Integral approximation.} For every integer weight $u$ with strictly decreasing
$u_\tau$ there are $C\ge0$, $s_n>0$ with $s_n\to0$, integers $q_n>0$ divisible by $8$, $\theta_n\in\Omega_M^\times$ and
$m_n\in\Z^{\mathcal R}$ with $m_n=q_n\,\mathrm{coef}(u_{s_n},s_n)$, such that, with
$L_n=L_{\varsigma(s_n)}(m_n)$: for all $n,\tau\in\N$ and $t\in\R$, every $\sigma\in\Lambda_\tau(y,a_{L_n})$ which is
small of degree $\tau$ for the shifted line $L_n(q_nt)$ satisfies
$\sigma\theta_n^\tau\in\deepstrmsmall\bigl(\mathrm{Or},\tau q_nu,\exp(\tau q_n(t+Cs_n))\bigr)$
(\lref{GlobalTheory.T.strict_weight_growth_bound}).
\end{enumerate}
\end{proposition}

Statement \ref{gi:descent} is the heart of the transport: a small section of the root model, divided by the gauge and
traced down to the function field of $Z$, is the value of a cusp form whose Hodge norms are controlled by the
degree of the flag and whose poles at the primes of $S(B)$ are bounded by the exponents $\deepstrex(p)$.

\subsubsection*{The arithmetic Hilbert--Samuel lower bound on root models}

\begin{lemma}\label{lem:pf:deepstr:shift}
\textup{(a)} For an intersection system $\deepstrII$, an arithmetic variety $P$ of dimension $D+1$, a smooth
Hermitian line $A$ and $t\in\R$, $\deepstrII(P;A(t)^{D+1})=\deepstrII(P;A^{D+1})+(D+1)\,t\,\deepstrgdeg(P,A)$.

\textup{(b)} Let $M$ be a root model, $\deepstrDim\ge1$, $\deepstrII$ an intersection system, $A$ a Hermitian line which
is nef on $P_M$ (semipositive relative to some $P_0\subseteq P_M$, nef on the fibres, conjugation-invariant),
$\kappa\in\R$, $r\ge1$, and $m\colon\N\to\N\cup\{\infty\}$ such that for every $\tau$ the number of $\sigma\in\Lambda_\tau(y,a_A)$
which are small of degree $\tau$ for $A(\kappa)$ is at most $m(8r\tau)$. If $\mathrm{(VI)}$ holds, then
\[
\max\Bigl(\frac{\deepstrII(P_M;A^{\deepstrDim+1})+(\deepstrDim+1)\kappa\,\deepstrgdeg(P_M,A)}
{(\deepstrDim+1)\,\delta_M\,(8r)^{\deepstrDim+1}},0\Bigr)\le\deepstrcvol(\deepstrDim,\delta_M,8;m).
\]
\formallabel{GlobalTheory.CompanionTheory.Ar.Small.selfInt_shiftH}
\formallabel{GlobalTheory.CompanionTheory.Ar.Plug.le_countVol_of_nef}
\formallabel{GlobalTheory.CompanionTheory.Ar.Plug.arithmeticVolume_div_le_countingVolume}
\end{lemma}

\begin{proof}
(a) $A(t)$ is $A\otimes\OO(-t)$ for the trivial line with constant metric; by the telescoping formula for
intersection systems (\lref{GroundTheory.Heights₂.Heights.Telescope.proof}) the difference of the two
self-intersections is a sum of $D+1$ integrals of the constant $t$ against products of Chern forms, each equal
to $t\,\deepstrgdeg(P,A)$ (\lref{GroundTheory.Heights₂.Heights.HilbertSamuel.TopConstantTwist.proof}).

(b) Write $D=\deepstrDim$. The shifted line $A(\kappa)$ is again nef
(\lref{GlobalTheory.T.hermitian_line_shift_properties}), so $\mathrm{(VI)}$ and (a) give
$\deepstrII(P_M;A^{D+1})+(D+1)\kappa\,\deepstrgdeg(P_M,A)\le\deepstrarvol(A(\kappa))$. Let $n_\tau$ be the number of small
$\sigma$ of degree $\tau$ ($n_\tau\ge1$ as $0$ is small). For $k=r\tau$,
\[
\frac{\log n_\tau}{\tau^{D+1}/(D+1)!}\cdot\frac1{(D+1)\,\delta_M\,(8r)^{D+1}}
=\frac{D!\,\log n_\tau}{\delta_M\,(8k)^{D+1}}\le\frac{D!\,\deepstrlogc(m(8k))}{\delta_M\,(8k)^{D+1}},
\]
and the $\limsup$ over the subsequence $k=r\tau$ is at most $\deepstrcvol(D,\delta_M,8;m)$. Hence
$\deepstrarvol(A(\kappa))/((D+1)\delta_M(8r)^{D+1})\le\deepstrcvol(D,\delta_M,8;m)$.
\end{proof}

\begin{lemma}[Limits along integral approximations]\label{lem:pf:deepstr:towerseq}
Let $T$ be a root tower, $\deepstrII$ a full intersection system with $\mathrm{(LC)}$, $(\lambda,g)$ with
$\mathrm{WH}(j,\lambda,g)$, $B$ with $\deepstrbnd_0(\lambda,g)<B+1$, $c\in\deepstrComp(T)$, $M=M(B,c)$, and let
$s_n>0$, $s_n\to0$, $q_n\ge1$ and $m_n=q_n\,\mathrm{coef}(\lambda_{s_n},s_n)\in\Z^{\mathcal R}$. Put
$L_n=L_{\varsigma(s_n)}(m_n)$ and $D=\deepstrDim$. Then $L_n$ is nef on $P_M$ for all large $n$, and as
$n\to\infty$
\[
\frac{\deepstrgdeg(P_M,L_n)}{q_n^{D}}\to\delta_M\,\deepstrvol_T(c)\,\deepstrflagdeg(\lambda),\qquad
\frac{\deepstrII(P_M;L_n^{D+1})}{(D+1)\,\deepstrgdeg(P_M,L_n)\,q_n}\to\deepstrslope_T(B,c,\lambda);
\]
moreover $\deepstrvol_T(c)\deepstrflagdeg(\lambda)>0$.
\formallabel{GlobalTheory.T.weight_threshold_limit_coefficients}
\formallabel{GlobalTheory.CompanionTheory.Ar.Small.tower_seq}
\end{lemma}

\begin{proof}
For all $n$, \cref{prop:pf:deepstr:combroot} gives
$\deepstrgdeg(P_M,L_n)=\delta_Mq_n^{D}\,m(B,c,\lambda,s_n)$ and
$\deepstrII(P_M;L_n^{D+1})=\delta_Mq_n^{D+1}\,\overline m(B,c,\lambda,s_n)$, which express intersection numbers of
integral combinations of the Chern-root lines through the functionals of \cref{def:pf:deepstr:rootmodel}. By
the limit conditions $\mathrm{(LC)}$, for large $n$ the model is nef on the fibres and semipositive at
$(\lambda_{s_n},s_n)$; then $L_n$ is nef on the fibres, semipositive and conjugation-invariant by
\cref{lem:pf:deepstr:nefcomb}, i.e.\ nef. The limits follow from those of $\mathrm{(LC)}$ along $s_n\to0^+$: the
first is $\delta_M\,m(B,c,\lambda,s_n)\to\delta_M\deepstrvol_T(c)\deepstrflagdeg(\lambda)$, and the second quotient equals
$\overline m/((D+1)m)$ at $s_n$ (with the convention $x/0=0$ of the development, both sides vanish if
$m(B,c,\lambda,s_n)=0$).
\end{proof}

\begin{proposition}[Volume of small model sections]\label{prop:pf:deepstr:volslope}
Assume $\sfname{Given}$ and $\mathrm{(VI)}$. In the situation of \cref{prop:pf:deepstr:geominputs}, let $u$ be an
integer weight, $\lambda=u$, $(g,B)$ with $\mathrm{WH}(j,\lambda,g)$ and $\deepstrbnd_0(\lambda,g)<B+1$, and
$b\in\deepstrComp(T)$. Then for every $t\in\R$
\begin{multline*}
\max\bigl(\deepstrvol_T(b)\,\deepstrflagdeg(\lambda)\,(\deepstrslope_T(B,b,\lambda)+t),0\bigr)\\
\le\deepstrcvol\bigl(\deepstrDim,\delta_{M(B,b)},8;\ \tau\mapsto\#\deepstrmsmall(\mathrm{Or}(B,b),\tau u,e^{\tau t})\bigr).
\end{multline*}
\formallabel{GlobalTheory.T.full_intersection_weight_bound}
\formallabel{GlobalTheory.CompanionTheory.Ar.Plug.volume_ge_slope}
\end{proposition}

\begin{proof}
Write $D=\deepstrDim$, $M=M(B,b)$, $\delta=\delta_M$, $V=\deepstrvol_T(b)\deepstrflagdeg(\lambda)$ and
$m(\tau)=\#\deepstrmsmall(\mathrm{Or},\tau u,e^{\tau t})$. Since $\lambda$ is regular dominant, each $u_\tau$ is strictly
decreasing. Take $C,s_n,q_n,\theta_n,m_n$ from \cref{prop:pf:deepstr:geominputs}\ref{gi:approx}, $L_n$ as there, $r_n=q_n/8$
and $t_n=t-Cs_n$. Note $D\ge1$: $\Psi_H\ne\varnothing$ and $p+q\ge2$ as $1\le j<p+q$. By \cref{lem:pf:deepstr:towerseq},
$L_n$ is nef for large $n$, and $G_n=\deepstrgdeg(P_M,L_n)>0$ for large $n$, as $G_n/q_n^D\to\delta V>0$.

\emph{Counting.} For every $\tau$, the map $\sigma\mapsto\sigma\theta_n^\tau$ sends the elements of $\Lambda_\tau(y,a_{L_n})$
which are small of degree $\tau$ for $L_n(q_nt_n)$ into $\deepstrmsmall(\mathrm{Or},\tau q_nu,\exp(\tau q_n(t_n+Cs_n)))
=\deepstrmsmall(\mathrm{Or},(q_n\tau)u,e^{q_n\tau t})$, whose cardinality is $m(q_n\tau)=m(8r_n\tau)$. By
\cref{lem:pf:deepstr:counting}(b), the hypothesis of \cref{lem:pf:deepstr:shift}(b) holds for $A=L_n$,
$\kappa=q_nt_n$, $r=r_n$. Hence, for large $n$, with $I_n=\deepstrII(P_M;L_n^{D+1})$,
\[
\max\Bigl(\frac{I_n+(D+1)q_nt_nG_n}{(D+1)\,\delta\,q_n^{D+1}},0\Bigr)\le\deepstrcvol(D,\delta,8;m).
\]
\emph{Limit.} With $h_n=I_n/((D+1)G_n)$, a direct computation gives
\[
\frac{I_n+(D+1)q_nt_nG_n}{(D+1)\,\delta\,q_n^{D+1}}=\frac1\delta\Bigl(\frac{h_n}{q_n}\cdot\frac{G_n}{q_n^{D}}
+t_n\frac{G_n}{q_n^{D}}\Bigr)
\]
by a direct computation. By \cref{lem:pf:deepstr:towerseq} and $t_n\to t$ the
right-hand side tends to $\frac1\delta(\deepstrslope_T(B,b,\lambda)\,\delta V+t\,\delta V)=V(\deepstrslope_T(B,b,\lambda)+t)$.
Since $x\mapsto\max(x,0)$ is continuous and the inequality holds for all large $n$, it holds in the limit.
\end{proof}

\subsubsection*{Transport to cusp forms}

\begin{proposition}[Transport of small sections]\label{prop:pf:deepstr:transport}
Assume $\sfname{Given}$, $\mathrm{(UP)}$ and $\mathrm{(VC)}$, in the situation of
\cref{prop:pf:deepstr:geominputs}. Let $Z$ be a component, $B_0\in\N$, $b=\kappa(Z)$, $\mathrm{Or}=\mathrm{Or}(B_0,b)$,
$(\mathrm{Fr},\mathrm{Rd})$ as in
\ref{gi:frame}, $u$ an integer weight with strictly decreasing $u_\tau$, $B_d>0$ with
$\deepstrPrimes(l,B_d)=S(B_0)$, $d=\deg\mathrm{Fl}>0$, and $k_0\in8\N$, $R_0\ge0$ and $f_1,\dots,f_d$ as in
\ref{gi:basis}. Let
$\tau\ge1$, $R\ge0$, $R'\in\R$ with $\deepstrclear(S(B_0),(\tau+k_0)u)\,d\,RR_0\le R'$. Then
\[
\#\deepstrmsmall(\mathrm{Or},\tau u,R)\le\bigl(\#\deepstrsmall_{B_d}(Z,(\tau+k_0)G_u,R')\bigr)^d.
\]
\formallabel{GlobalTheory.CompanionTheory.Ar.Plug.SlopeVol.modelSmall_encard_le}
\end{proposition}

\begin{proof}
Put $S=S(B_0)$, $u'=(\tau+k_0)u$ (so $(\tau+k_0)G_u=G_{u'}$ and $u'_\tau$ is strictly decreasing) and
$N=\deepstrclear(S,u')\ge1$. Consider the relation $\mathsf P(s,i,f')$: \emph{$f'=Nf$ for a rational section
$f$ with $\deepstrev(f)=\deepstrtr_{u'}(sf_i)$}. We check the hypotheses of \cref{lem:pf:deepstr:counting}(a)
with $\mathcal A=\deepstrmsmall(\mathrm{Or},\tau u,R)$ and $\mathcal B=\deepstrsmall_{B_d}(Z,G_{u'},R')$.

\emph{Existence.} Let $s\in\mathcal A$ and $1\le i\le d$. By \ref{gi:mul}, $sf_i\in\deepstrmsmall(\mathrm{Or},u',RR_0)$. By
\ref{gi:descent} there is $f\in\deepstrcusp(Z,G_{u'})$ with $\deepstrev(f)=\deepstrtr_{u'}(sf_i)$,
$\|f\|_{G_{u'},\sigma,z}\le dRR_0$ and $v(f)\ge-v_\lambda(p)[p\in S]\deepstrex(p)\|u'\|_1$ for $v\in\mathcal
V^{\mathrm{rec}}(\lambda)$. By \cref{lem:pf:deepstr:clear}(b) (with $S=\deepstrPrimes(l,B_d)$), $w(Nf)\ge0$ for all
$w\in\mathcal V_{B_d}(\lambda)$, and $\|Nf\|_{G_{u'},\sigma,z}=N\|f\|_{G_{u'},\sigma,z}\le NdRR_0\le R'$. So
$f'=Nf\in\mathcal B$ and $\mathsf P(s,i,f')$.

\emph{Injectivity.} Let $s,s'\in\mathcal A$ and suppose that for every $i$ some $f'$ satisfies $\mathsf P(s,i,f')$
and $\mathsf P(s',i,f')$: $f'=Nf=Ng$ with $\deepstrev(f)=\deepstrtr_{u'}(sf_i)$ and $\deepstrev(g)=\deepstrtr_{u'}(s'f_i)$.
As $N\ne0$, $f=g$, so $\deepstrtr_{u'}((s-s')f_i)=0$ for all $i$ (the trace is additive), and $s=s'$ by
\ref{gi:basis}. \Cref{lem:pf:deepstr:counting}(a) gives the claim.
\end{proof}

\subsubsection*{The comparison theorem}

\begin{theorem}[Comparison of slopes]\label{thm:pf:deepstr:comparison}
Assume \ref{H:pf:deepstr:nef}, $\sfname{Given}$ and $\mathrm{(VI)}$. Let $E$ be a CM field Galois over $\Q$ with
host $H$; $D_w,l,X^o,c$; $r\in\N$ and uniformisations $U$ with the composition laws $\mathrm{(CL)}$; assume $\mathrm{(UP)}$;
let $x$ with $|\tau(\xi_0)|=x$ for all good $\sigma$; let $Y$ be a model of $X^o$ over $\OO_{F_c}$, $T$ a root tower of
type $(\Psi_H,p+q)$ with a good origin $o$, $\kappa\colon\deepstrComp(X^o,c)\to\deepstrComp(T)$ a bijection compatible
with $o$, and $\deepstrII$ a full intersection system such that $T$ is proportional and satisfies
$\mathrm{(LC)}$. Let $Z$ be a component, $u$ an integer weight with strictly decreasing $u_\tau$, $\lambda=u$,
and $g,B\in\R$ with $\mathrm{WH}(j,\lambda,g)$ and $\deepstrbnd_0(\lambda,g)\le B$. Then
\[
\deepstrslope_T(\lfloor B\rfloor,\kappa(Z),\lambda)-2\deepstrcost(l,B)\,\|u\|_1\le\deepstrvsl(Z,u),
\]
the volume slope \eqref{eq:pf:deepstr:compvol} being computed with $\mathcal V_B$.
\formallabel{GlobalTheory.T.uniform_all_norm_height_bound}
\formallabel{GlobalTheory.CompanionTheory.Relay.uniform_all_norm_height_bound}
\formallabel{GlobalTheory.CompanionTheory.Ar.Plug.SlopeVolM.slope_sub_cost_le_compVolSlope}
\end{theorem}

\begin{proof}
Put $B_0=\lfloor B\rfloor$, $b=\kappa(Z)$, $\mathrm{Or}=\mathrm{Or}(B_0,b)$, $M=M(B_0,b)$, $D=\deepstrDim$, $n=[F_c:\Q]$,
$d=\deg\mathrm{Fl}$, $S=S(B_0)$, $L=\|u\|_1$ and $\sigma_0=\deepstrslope_T(B_0,b,\lambda)-2\deepstrcost(l,B)L$.

\emph{Preliminaries.} By \cref{prop:pf:deepstr:geominputs}\ref{gi:topdeg}, $\deepstrtop(Z,u)=\deepstrvol_T(b)\deepstrflagdeg(\lambda)>0$.
As $\deepstrbnd_0(\lambda,g)\le B<B_0+1$, \ref{gi:frame} and \ref{gi:basis} provide $(\mathrm{Fr},\mathrm{Rd})$ and $k_0\in8\N$,
$R_0>0$, $f_1,\dots,f_d$; \ref{gi:htlow} gives $\mathrm{(VC)}$. We have $\deepstrPrimes(l,B)=S$: a prime $p$
satisfies $p\le B$ if and only if $p\le B_0$, and $p\in\mathrm{bad}(T)$ if and only if $p\mid2d_EM'$, i.e.\ $p\mid2|d_E|M'$.
Consequently $\deepstrcost(l,B)=\sum_{p\in S}\deepstrex(p)\log p$. By \eqref{eq:pf:deepstr:deltaeq}, $\delta_M=nd$, and
$\delta_M>0$ gives $d>0$.

\emph{Reduction.} By \eqref{eq:pf:deepstr:volslopeiff} we must show, for every $t\in\R$,
$\max(\deepstrtop(Z,u)(\sigma_0+t),0)\le\mathrm{cv}_{Z,u}(t)$, with $\mathrm{cv}_{Z,u}$ as in
\cref{sec:pf:deepstr:conv}. By continuity it suffices to show
$\max(\deepstrtop(Z,u)(\sigma_0+t-\eps),0)\le\mathrm{cv}_{Z,u}(t)$ for every $\eps>0$. Fix $t$ and $\eps>0$, put
$t'=t-2\deepstrcost(l,B)L-\eps$, and note
\[
\deepstrtop(Z,u)(\sigma_0+t-\eps)=\deepstrvol_T(b)\deepstrflagdeg(\lambda)\bigl(\deepstrslope_T(B_0,b,\lambda)+t'\bigr).
\]
By \cref{prop:pf:deepstr:volslope} (at $B_0$, $b$ and $t'$), the positive part of this number is at most
$\deepstrcvol(D,nd,8;m)$ with $m(\tau)=\#\deepstrmsmall(\mathrm{Or},\tau u,e^{\tau t'})$. Let
$m'(\tau)=\#\deepstrsmall_B(Z,\tau G_u,e^{\tau t})$; then $\mathrm{cv}_{Z,u}(t)=\deepstrcvol(D,n,8;m')$, and by
\cref{lem:pf:deepstr:counting}(c) it suffices to show $m(\tau)\le m'(\tau+k_0)^d$ for all large $\tau$.

\emph{Transport.} Let $c_S=\sum_{p\in S}\log p$. Choose $\tau\ge1$ so large that
$2k_0L\deepstrcost(l,B)+c_S+\log(dR_0)-k_0t\le\tau\eps$. By \cref{lem:pf:deepstr:clear}(a) and $\|(\tau+k_0)u\|_1=(\tau+k_0)L$,
\[
\deepstrclear(S,(\tau+k_0)u)\,d\,e^{\tau t'}R_0\le\exp\bigl(2(\tau+k_0)L\deepstrcost(l,B)+c_S+\tau t'+\log(dR_0)\bigr)
\le e^{(\tau+k_0)t},
\]
because $2(\tau+k_0)L\deepstrcost(l,B)+\tau t'=\tau t-\tau\eps+2k_0L\deepstrcost(l,B)$. Hence
\cref{prop:pf:deepstr:transport} (with $B_d=B$, $R=e^{\tau t'}$, $R'=e^{(\tau+k_0)t}$) gives $m(\tau)\le m'(\tau+k_0)^d$.
\end{proof}

% ---------------------------------------------------------------------------
\subsection{Proofs of further results}\label{sec:pf:deepstr:further}

This subsection proves further results of the formal development which are used above and in other
sections of Part II: symmetry and multilinearity of geometric intersection numbers; constant metrics, fibre
degrees and expansions of arithmetic intersection numbers and Chern-form integrals of integral combinations of
lines on root models; nefness statements; facts on generic flags and flag varieties; the existence of an
integral section attached to an origin; and two Hodge--Tate statements. For smooth Hermitian lines $L_1,\dots,L_e$ on an arithmetic variety $P$ of dimension $e+1$ we
write $\int_Pc_1(L_1)\cdots c_1(L_e)$ for the integral over the complex points of $P$ of the product of their
Chern forms (the quantity used in \cref{def:pf:deepstr:rootmodel}). For lines $\ell=(\ell_j)_{j\in J}$ and
$n\in\Z^J$ put $\ell(n)=\bigotimes_j\ell_j^{\otimes n_j}$ (multidegree $\sum_jn_ja_{\ell_j}$, metric $\prod_jh_{\ell_j}^{n_j}$).
For an ideal $V\subseteq R$ and multidegrees $a_1,\dots,a_n\in\Z^K$ the \emph{geometric degree}
$\deg_V(a_1,\dots,a_n)\in\Z$ is defined recursively: it is $0$ if $V$ is not relevant; otherwise choose for each
block $k$ a variable $x_{k,i_k}\notin V$, let $\deg_V()$ be the length of the $\Z$-module $R/(V+(x_{k,i_k}-1)_k)$, and put
\[
\deg_V(a_1,\dots,a_{n+1})=\sum_ka_{n+1,k}\sum_W\mathrm{mult}_W\bigl(V+(x_{k,i_k})\bigr)\deg_W(a_1,\dots,a_n),
\]
the inner sum over the relevant primes $W$ minimal over $V+(x_{k,i_k})$.

\subsubsection*{Geometric intersection numbers}

\begin{definition}[Field intersection families]\label{def:pf:deepstr:fieldint}
Let $F$ be a field and $F[x]=F[x_{k,i}]$ the multigraded polynomial ring of \cref{sec:pf:deepstr:conv}. A
multihomogeneous prime $V\subset F[x]$ is a \emph{projective variety of dimension $d$} if it is relevant and
$\dim F[x]/V=d+|K|$. A \emph{rational section} of multidegree $a\in\Z^K$ on $V$ is a pair of multihomogeneous
$f,g\notin V$ with $\deg f-\deg g=a$; its divisor $\mathrm{div}_V(f/g)=\mathrm{div}_V(f)-\mathrm{div}_V(g)$ is a finite
formal sum of primes $W$, $\mathrm{div}_V(f)$ being supported on the relevant primes minimal over $V+(f)$. A
\emph{field intersection family} assigns an integer $I(V;L_1,\dots,L_n)$ to every
multihomogeneous ideal $V$ and $L_1,\dots,L_n\in\Z^K$ such that: (i) for $V$ of dimension $0$, $I(V;\varnothing)$
is the degree of the point $V$; (ii) for $V$ of dimension $n+1$ and every rational section $f/g$ of
multidegree $A$ on $V$,
\[
I(V;L_1,\dots,L_n,A)=\sum_W\mathrm{ord}_W(f/g)\,I(W;L_1,\dots,L_n).
\]
This is the definition \lref{ReciprocityTheory.ChartRestriction.G.AI.I1.IsFieldIntersection}.
\end{definition}

\begin{lemma}[Standard sections and divisors]\label{lem:pf:deepstr:stdsec}
Let $V\subset F[x]$ be a relevant multihomogeneous prime.
\begin{enumerate}[label=(\alph*)]
\item For every $a\in\Z^K$ there is a rational section $f/g$ of multidegree $a$ on $V$ with $f$ and $g$ monomials. If
$V$ is a projective variety of dimension $n+1$, then for every field intersection family $I$ and
$L_1,\dots,L_n\in\Z^K$ we have $I(V;L_1,\dots,L_n,a)=\sum_W\mathrm{ord}_W(f/g)\,I(W;L_1,\dots,L_n)$.
\item Let $V$ be a projective variety of dimension $n+1$. If $f\notin V$ is multihomogeneous, every relevant prime
$W$ minimal over $V+(f)$ is a projective variety of dimension $n$. Consequently, if $f/g$ is a rational section on
$V$ and $\mathrm{ord}_W(f/g)\ne0$, then $W$ is a projective variety of dimension $n$.
\end{enumerate}
\formallabel{ReciprocityTheory.ChartRestriction.G.AI.I1.isRatSection_sec}
\formallabel{ReciprocityTheory.ChartRestriction.G.AI.I1.multi_eval_sec}
\formallabel{ReciprocityTheory.ChartRestriction.G.AI.I1.isProjVarF_of_mem_divSupp}
\formallabel{ReciprocityTheory.ChartRestriction.G.AI.I1.mem_of_ratDiv_ne_zero}
\formallabel{ReciprocityTheory.ChartRestriction.G.AI.I1.multi_isProjVarF_of_ratDiv}
\end{lemma}

\begin{proof}
(a) As $V$ is relevant, for each block $k$ there is a variable $x_{k,i_k}\notin V$. Put
$f=\prod_kx_{k,i_k}^{\max(a_k,0)}$ and $g=\prod_kx_{k,i_k}^{\max(-a_k,0)}$; they are multihomogeneous, not in $V$ because $V$
is prime, and $\deg f-\deg g=a$. The formula is (ii) for this section.
(b) The minimal primes over the multihomogeneous ideal $V+(f)$ are multihomogeneous
(\lref{ReciprocityTheory.ChartRestriction.G.AI.I1.cyc_isMultiHomogIdeal_of_mem_minimalPrimes}). By Krull's
principal ideal theorem in the domain $F[x]/V$, which is finitely generated over $F$, every prime $W$ minimal over
$V+(f)$ satisfies $\dim F[x]/W+1=\dim F[x]/V$
(\lref{ReciprocityTheory.ChartRestriction.G.AI.I1.cyc_dim_quot_minimal_add_one}); so $\dim F[x]/W=n+|K|$, and $W$
is a projective variety of dimension $n$ if it is relevant. If $\mathrm{ord}_W(f/g)=\mathrm{ord}_W(f)-\mathrm{ord}_W(g)\ne0$, then $W$
lies in the support of $\mathrm{div}_V(f)$ or of $\mathrm{div}_V(g)$, that is, it is a relevant prime minimal over $V+(f)$ or
over $V+(g)$, where $f,g\notin V$ are multihomogeneous.
\end{proof}

\begin{theorem}[Symmetry]\label{thm:pf:deepstr:fieldsym}
Let $I$ be a field intersection family, $V$ a projective variety of dimension $n$, $L\in(\Z^K)^n$ and $\sigma$ a
permutation of $\{1,\dots,n\}$. Then $I(V;L_{\sigma(1)},\dots,L_{\sigma(n)})=I(V;L_1,\dots,L_n)$.
\formallabel{ReciprocityTheory.ChartRestriction.G.AI.I1.intersection_symmetric}
\end{theorem}

\begin{proof}
Induction on $n$; the case $n=0$ is trivial. Let $V$ have dimension $n+1$. (a) \emph{Transpositions of two
of the first $n$ entries.} By \cref{lem:pf:deepstr:stdsec}(a), applied to a rational section of multidegree
$L_{n+1}$, $I(V;L)=\sum_W\mathrm{ord}_W\,I(W;L_1,\dots,L_n)$, and every $W$ with nonzero coefficient is a projective
variety of dimension $n$ by \cref{lem:pf:deepstr:stdsec}(b); by induction each term is unchanged.
(b) \emph{The transposition of the last two entries}: $I(V;L',a,b)=I(V;L',b,a)$ for every projective variety
$V$ of dimension $m+2$ is \lref{ReciprocityTheory.ChartRestriction.G.AI.I1.multi_swap_last_two} (the
commutativity of the intersection of two divisors, compare \cite[Thm.~2.4]{Fulton1998}). (c) A transposition
$(x,n+1)$ with $x<n$ equals $(n,n+1)(x,n)(n,n+1)$; apply (b), (a) and (b). (d) Every permutation is a
product of transpositions.
\end{proof}

\begin{theorem}[Multilinearity]\label{thm:pf:deepstr:fieldmultilin}
Let $I$ be a field intersection family, $V$ a projective variety of dimension $n$, $L\in(\Z^K)^n$, $1\le j\le n$ and
$a,b\in\Z^K$. Then the value of $I(V;L)$ with $L_j$ replaced by $a+b$ is the sum of the values with $L_j$ replaced
by $a$ and by $b$.
\formallabel{ReciprocityTheory.ChartRestriction.G.AI.I1.multilin_thm}
\formallabel{ReciprocityTheory.ChartRestriction.G.AI.I1.intersection_multilinear}
\formallabel{ReciprocityTheory.ChartRestriction.G.AI.I1.multilin_last}
\end{theorem}

\begin{proof}
First let $j=n$. If $f/g$ and $f'/g'$ are rational sections of multidegrees $a$ and $b$, then $ff'/gg'$ is a
rational section of multidegree $a+b$ and $\mathrm{div}_V(ff'/gg')=\mathrm{div}_V(f/g)+\mathrm{div}_V(f'/g')$, the orders
being valuations; so (ii) and the linearity of $D\mapsto\sum_W\mathrm{ord}_W(D)\,I(W;L_1,\dots,L_{n-1})$ give the claim.
For general $j$, move the $j$-th entry to the last position by the transposition $(j,n)$, which does not
change the values by \cref{thm:pf:deepstr:fieldsym}.
\end{proof}

\subsubsection*{Arithmetic intersection numbers of combinations}

\begin{lemma}[One slot]\label{lem:pf:deepstr:oneslot}
Let $\deepstrII$ be a full intersection system, $P$ an arithmetic variety of dimension $d$, $L_1,\dots,L_d$
smooth Hermitian lines, $1\le k\le d$, $\ell=(\ell_j)_{j\in J}$ smooth lines and $n\in\Z^J$. Then
$\deepstrII(P;L)$ with $L_k$ replaced by $\ell(n)$ equals $\sum_jn_j\,\deepstrII(P;L)$ with $L_k$ replaced by $\ell_j$.
\formallabel{RootFieldNumerics.intersection_multilinear_one_position}
\end{lemma}

\begin{proof}
For a full system, $\varphi(A)=\deepstrII(P;L)$ with $L_k$ replaced by $A$ satisfies $\varphi(A\otimes B)=\varphi(A)+\varphi(B)$
for smooth $A,B$ (multilinearity, part of the definition of a full system). The lines $\ell(m+m')$ and
$\ell(m)\otimes\ell(m')$ have the same multidegree and the same metric, so $m\mapsto\varphi(\ell(m))$ is a homomorphism
$\Z^J\to\R$; as $\ell(e_j)=\ell_j$ for the basis vectors, $\varphi(\ell(n))=\sum_jn_j\varphi(\ell_j)$.
\end{proof}

\begin{corollary}[Expansion]\label{cor:pf:deepstr:expansion}
In the situation of \cref{lem:pf:deepstr:oneslot}, for $n_1,\dots,n_d\in\Z^J$,
\[
\deepstrII\bigl(P;\ell(n_1),\dots,\ell(n_d)\bigr)=\sum_{w\in J^d}\Bigl(\prod_{k=1}^dn_k(w_k)\Bigr)\deepstrII(P;\ell_{w_1},\dots,\ell_{w_d}).
\]
\formallabel{RootFieldNumerics.intersection_multilinear_expansion}
\formallabel{GlobalTheory.T.arith_intersection_multilinear_all}
\end{corollary}

\begin{proof}
Apply \cref{lem:pf:deepstr:oneslot} successively in the slots $1,\dots,d$ (all lines $\ell(m)$ are smooth) and
expand the product of the sums.
\end{proof}

\begin{lemma}[Constant metrics and fibre degrees]\label{lem:pf:deepstr:constmetric}
Let $\deepstrII$ be a full intersection system, $P$ an arithmetic variety of dimension $e+1$ and $L_1,\dots,L_e$ smooth
Hermitian lines. For $c\in\R$ let $O_c$ be the line of multidegree $0$ with the constant metric $h\equiv e^{-c}$.
\begin{enumerate}[label=(\alph*)]
\item $\deepstrII(P;L_1,\dots,L_e,O_c)=-c\int_Pc_1(L_1)\cdots c_1(L_e)$.
\item For every prime $p$,
\[
\int_Pc_1(L_1)\cdots c_1(L_e)=\sum_V\mathrm{mult}_V(P+(p))\,\deg_V(\deg L_1,\dots,\deg L_e),
\]
the sum over the relevant primes $V$ minimal over $P+(p)$.
\end{enumerate}
\formallabel{GlobalTheory.T.arith_intersection_constant_metric}
\formallabel{GlobalTheory.CompanionTheory.Ar.Top.intersection_snoc_constWeight}
\formallabel{RootFieldNumerics.intersection_constant_metric_line}
\formallabel{GlobalTheory.T.intersection_system_degree}
\formallabel{GlobalTheory.CompanionTheory.Ar.Top.degreeStatement_of_intersectionSystem}
\end{lemma}

\begin{proof}
(a) The line $O_c$ is smooth, its metric being constant. Apply the recursion of \cref{def:covolume:hs}(c) to the
last line $O_c$ and its section $s=1$, which is multihomogeneous of multidegree $0$ and not in $P$. The ideal $P+(1)$
is the unit ideal, so there are no components $W$, and $\log\|1\|=\log(1/h)=c$ identically; hence
$\deepstrII(P;L_1,\dots,L_e,O_c)=-\int_Pc\,c_1(L_1)\cdots c_1(L_e)$.
(b) Apply the recursion to the line $O_0$ (trivial metric) and its section $p$, which is multihomogeneous of
multidegree $0$ and not in $P$ because $P\cap\Z=0$. Every component $V$ of $P+(p)$ lies over $p$ and contributes
$\mathrm{mult}_V(P+(p))\deg_V(\deg L_1,\dots,\deg L_e)\log p$, and $\log\|p\|=\log p$ identically. By (a) with $c=0$ the
left-hand side of the recursion vanishes, so
\[
0=\log p\Bigl(\sum_V\mathrm{mult}_V(P+(p))\deg_V(\deg L_1,\dots,\deg L_e)-\int_Pc_1(L_1)\cdots c_1(L_e)\Bigr).
\]
In the development, (b) is obtained from the conservation of degrees in the fibres for full intersection systems
(\lref{GlobalTheory.CompanionTheory.Ar.Giv.degree_of_isIntersectionSystem}, a result of the shared part
stating (b) for every full intersection system).
\end{proof}

\begin{lemma}[Chern-form integrals]\label{lem:pf:deepstr:chernint}
Let $\deepstrII$ be a full intersection system, $P$ an arithmetic variety of dimension $e+1$ and
$\ell=(\ell_j)_{j\in J}$ smooth lines.
\begin{enumerate}[label=(\alph*)]
\item For $n_1,\dots,n_e\in\Z^J$,
\[
\int_Pc_1(\ell(n_1))\cdots c_1(\ell(n_e))=\sum_{w\in J^e}\Bigl(\prod_kn_k(w_k)\Bigr)\int_Pc_1(\ell_{w_1})\cdots c_1(\ell_{w_e}).
\]
\item If $L_1,\dots,L_e$ and $L'_1,\dots,L'_e$ are smooth lines and $\sigma$ is a permutation with
$\deg L'_k=\deg L_{\sigma(k)}$ for all $k$, then $\int_Pc_1(L'_1)\cdots c_1(L'_e)=\int_Pc_1(L_1)\cdots c_1(L_e)$.
\end{enumerate}
\formallabel{GlobalTheory.T.log_integral_multilinear}
\formallabel{GlobalTheory.CompanionTheory.Ar.Top.geo_multilin}
\formallabel{GlobalTheory.CompanionTheory.Ar.Top.GeoMix.pts_eq_II}
\formallabel{GlobalTheory.CompanionTheory.Ar.Top.GeoMix.pts_comb}
\formallabel{GlobalTheory.CompanionTheory.Ar.Top.GeoMix.pts_perm}
\formallabel{GlobalTheory.CompanionTheory.Ar.Top.GeoMix.pts_deg}
\end{lemma}

\begin{proof}
With $O_c$ as in \cref{lem:pf:deepstr:constmetric}, part (a) of that lemma with $c=-1$ gives
\begin{equation}\label{eq:pf:deepstr:ptsII}
\int_Pc_1(L_1)\cdots c_1(L_e)=\deepstrII(P;L_1,\dots,L_e,O_{-1}).
\end{equation}
(a) Extend $\ell$ by $\ell_\ast=O_{-1}$ and the exponent vectors by $0$ at $\ast$, and let the last slot carry the
exponent vector $e_\ast$. By \eqref{eq:pf:deepstr:ptsII} and \cref{cor:pf:deepstr:expansion}, the left-hand side is
a sum over words $v$ of length $e+1$ in $J\sqcup\{\ast\}$; the coefficient $\prod_kn_k(v_k)\cdot e_\ast(v_{e+1})$ vanishes
unless $v_{e+1}=\ast$ and $v_k\in J$ for $k\le e$, and the remaining terms are, by \eqref{eq:pf:deepstr:ptsII} again,
those of the right-hand side.
(b) By \cref{lem:pf:deepstr:constmetric}(b) the integral depends only on the multidegrees; so
$\int_P\prod_kc_1(L'_k)=\int_P\prod_kc_1(L_{\sigma(k)})$. The latter equals $\int_P\prod_kc_1(L_k)$ by
\eqref{eq:pf:deepstr:ptsII} and the symmetry of intersection systems under permutations of the lines (the
permutation fixing the last slot, which carries $O_{-1}$).
\end{proof}

\begin{proposition}[Combinations on root models]\label{prop:pf:deepstr:combroot}
Let $\deepstrII$ be a full intersection system, $M$ a root model with $D=\deepstrDim$, $t\in\R$ and $q\in\N$.
\begin{enumerate}[label=(\alph*)]
\item If $\lambda$ is a weight, $s\in\R$ and $n=q\,\mathrm{coef}(\lambda,s)\in\Z^{\mathcal R}$, then
$\deepstrII(P_M;L_t(n)^{D+1})=\delta_Mq^{D+1}\,\mathrm{ar}_t\bigl(\ell_s(\lambda)^{D+1}\bigr)$, with $\ell_s(\lambda)=\deepstrlin(\lambda)+sX_\ast$.
\item If $\lambda_1,\dots,\lambda_D$ are weights, $s_1,\dots,s_D\in\R$ and $n_k=q\,\mathrm{coef}(\lambda_k,s_k)\in\Z^{\mathcal R}$,
then $\int_{P_M}c_1(L_t(n_1))\cdots c_1(L_t(n_D))=\delta_Mq^D\,\mathrm{geo}\bigl(\prod_k\ell_{s_k}(\lambda_k)\bigr)$.
\end{enumerate}
\formallabel{GlobalTheory.CompanionTheory.Ar.Top.top_comb_eq_arFun}
\formallabel{RootFieldNumerics.root_model_arithmetic_intersection}
\formallabel{GlobalTheory.CompanionTheory.Ar.Top.geoFun_comb_mixed}
\formallabel{RootFieldNumerics.root_model_points_integral}
\end{proposition}

\begin{proof}
(a) By \cref{cor:pf:deepstr:expansion}, $\deepstrII(P_M;L_t(n)^{D+1})=\sum_w\prod_kn(w_k)\,\deepstrII(P_M;x_t(w_1),\dots,x_t(w_{D+1}))$,
the sum over words $w\in\mathcal R^{D+1}$. Write $\ell_s(\lambda)=\sum_rc_rX_r$ with $c_r=-\lambda_r$ for $r\ne\ast$ and
$c_\ast=s$, so that $c_r=\mathrm{coef}(\lambda,s)_r$ for $r\ne\ast$ and $c_\ast=k_A\,\mathrm{coef}(\lambda,s)_\ast$. Expanding,
$\ell_s(\lambda)^{D+1}=\sum_w\prod_kc_{w_k}X^{m(w)}$, where $m(w)_r$ is the number of $k$ with $w_k=r$. The lines of
the monomial $X^{m(w)}$, listed in the order used in \cref{def:pf:deepstr:rootmodel}, are a permutation of
$x_t(w_1),\dots,x_t(w_{D+1})$, so by the symmetry of intersection systems
$\mathrm{ar}_t(X^{m(w)})=\deepstrII(P_M;x_t(w_1),\dots,x_t(w_{D+1}))/(\delta_Mk_A^{m(w)_\ast})$. Since
$\prod_kc_{w_k}=k_A^{m(w)_\ast}\prod_k\mathrm{coef}(\lambda,s)_{w_k}$ and $\prod_kn(w_k)=q^{D+1}\prod_k\mathrm{coef}(\lambda,s)_{w_k}$,
\[
\delta_Mq^{D+1}\,\mathrm{ar}_t(\ell_s(\lambda)^{D+1})=\sum_wq^{D+1}\prod_k\mathrm{coef}(\lambda,s)_{w_k}\,
\deepstrII(P_M;x_t(w_1),\dots)=\deepstrII(P_M;L_t(n)^{D+1}).
\]
(b) The same computation with \cref{lem:pf:deepstr:chernint}(a) in place of \cref{cor:pf:deepstr:expansion};
for each word $w\in\mathcal R^D$, \cref{lem:pf:deepstr:chernint}(b) replaces the lines $x_t(w_k)$ by the lines
$x_1(r)$ of the monomial $X^{m(w)}$ (they have the same multidegrees up to the permutation), whence
$\int_{P_M}\prod_kc_1(x_t(w_k))=\delta_Mk_A^{m(w)_\ast}\,\mathrm{geo}(X^{m(w)})$.
\end{proof}

\begin{lemma}[Nefness and semipositivity of combinations]\label{lem:pf:deepstr:nefcomb}
Let $M$ be a root model, $t\in\R$, $\lambda$ a weight, $s\in\R$, $q\in\N$ and $n=q\,\mathrm{coef}(\lambda,s)\in\Z^{\mathcal R}$. For an
ideal $W$ let $\deg_W$ be the geometric degree with one argument, and put
$\mathrm{cd}(\lambda,s,W)=\sum_r\mathrm{coef}(\lambda,s)_r\deg_W(a_r)$.
\begin{enumerate}[label=(\alph*)]
\item $\deg_W$ is linear: $\deg_W(\sum_jm_jb_j)=\sum_jm_j\deg_W(b_j)$ for $m\in\Z^J$ and $b_j\in\Z^K$. In particular
$\deg_W(\deg L_t(n))=q\,\mathrm{cd}(\lambda,s,W)$ for every $W$; hence, if $\mathrm{cd}(\lambda,s,W)\ge0$ for every vertical curve
$W\supseteq P_M$ (\emph{$M$ is nef on the fibres at $(\lambda,s)$}), then $L_t(n)$ is nef on the fibres.
\item If the Levi form of $w_{t,\lambda,s}=\sum_r\mathrm{coef}(\lambda,s)_r\log h_{x_t(r)}$ is non-negative on the cone-tangent
vectors at the complex points of $P_M$ (\emph{$M$ is semipositive at $(t,\lambda,s)$}), then $L_t(n)$ is semipositive
on $P_M$ and conjugation-invariant.
\end{enumerate}
\formallabel{GlobalTheory.CompanionTheory.Ar.Top.nefFib_integral_multiple}
\formallabel{RootFieldNumerics.curveDeg_eq_and_nefOnFibres_of_coef}
\formallabel{RootFieldNumerics.root_model_line_semipositive}
\formallabel{RootFieldNumerics.geoNum_linear_combination}
\formallabel{RootFieldNumerics.geoNum_one_linear_in_degree}
\end{lemma}

\begin{proof}
(a) By the definition of the geometric degree with one argument, $\deg_W(b)=\sum_kb_k\,d_k(W)$, where
$d_k(W)=\sum_V\mathrm{mult}_V(W+(x_{k,i_k}))\deg_V()$ if $W$ is relevant and $d_k(W)=0$ otherwise; so $\deg_W$ is linear,
with $d_k(W)=\deg_W(e_k)$. The multidegree of $L_t(n)$ is $\sum_rn_ra_r$, independent of $t$; substitute
$n_r=q\,\mathrm{coef}(\lambda,s)_r$. A vertical curve $W\supseteq P_M$ has prime characteristic $p$, so the hypothesis applies
and $\deg_W(\deg L_t(n))=q\,\mathrm{cd}(\lambda,s,W)\ge0$. (b) $L_t(n)$ is smooth and conjugation-invariant as a
tensor product of such lines. Near every complex point of the multicone with non-vanishing blocks,
$\log h_{L_t(n)}=\sum_rn_r\log h_{x_t(r)}=q\,w_{t,\lambda,s}$, so its Levi form is $q$ times that of $w_{t,\lambda,s}$, hence
non-negative.
\end{proof}

\subsubsection*{Nefness on the generic fibre}

\begin{definition}[Analytic systems, realisations and real towers]\label{def:pf:deepstr:ansys}
(a) An \emph{analytic system} $S$ of dimension $D$ with index set $\sigma$ consists of a compact Hausdorff space
$X_S$ with a holomorphic atlas of dimension $D$, a positive potential, a line potential with a section whose
divisor has strict normal crossings, and for each $a\in\sigma$ a singular weight and a family of regularised
potentials $\mathrm{pot}(a,t)$, $0<t\le t_0$; it has a \emph{degree polynomial} $\deg_S\in\R[X_a:a\in\sigma]$.
(b) A \emph{realisation} $\mathrm{Re}$ of $S$ on an arithmetic variety $P$ of dimension $D+1$ consists of a
logarithmic resolution of $P$ along an ideal $J$, a holomorphic map $\pi$ from it to $X_S$ of degree
$\deg\mathrm{Re}>0$, smooth \emph{reference lines} $\mathrm{ref}(a)$ on $P$ whose potentials are the pull-backs of the
singular weights, and regularised lines with the same multidegrees. Put $\mathrm{ref}(n)=\bigotimes_a\mathrm{ref}(a)^{\otimes n_a}$.
(c) A \emph{real tower} $R$ over a root tower $T$ of type $(\Psi,N)$ consists of an analytic system $S_c$ with
index set $\mathcal R$ for each component $c$ and, for all $B$ and $c$, a realisation of $S_c$ on $P_{M(B,c)}$ whose
reference lines satisfy $\deg\mathrm{ref}(\tau,i)=-\sum_{k\le i}a_{(\tau,k)}$ and $\deg\mathrm{ref}(\ast)=\deg A$; $R$ is \emph{smooth} if $J=(1)$ and the generic fibres of the models are smooth, and \emph{flat} if the
potentials of the indices $(\tau,N-1)$ are pluriharmonic.
(d) A weight $\lambda$ is \emph{dominant} if $\lambda_{\tau,b}\le\lambda_{\tau,a}$ for $a\le b$. A root model $M$ is \emph{nef on the
generic fibre} at $(\lambda,s)$ if $\sum_r\mathrm{coef}(\lambda,s)_r\,\deepstrgdeg(W,x_1(r))\ge0$ for every arithmetic curve
$W\supseteq P_M$.
(e) The \emph{nef polarisation statement} is the statement of \cref{thm:pf:deepstr:nefpol} below for all
data.
\formallabel{GlobalTheory.CompanionTheory.Ar.Top.RealTower}
\formallabel{GlobalTheory.CompanionTheory.Ar.Top.Dominant}
\formallabel{GlobalTheory.CompanionTheory.Ar.PB.NefPol}
\end{definition}

\begin{lemma}[Chern-form integrals and the generic fibre]\label{lem:pf:deepstr:genericfibre}
Assume \ref{G:pf:deepstr:pb}. Let $I$ be a field intersection family over $\Q$, $P$ an arithmetic variety of dimension
$D+1$ with generic fibre $P_\Q=P\cdot\Q[x]$, and $L_1,\dots,L_D$ smooth Hermitian lines. Then
$\int_Pc_1(L_1)\cdots c_1(L_D)=I(P_\Q;\deg L_1,\dots,\deg L_D)$.
\formallabel{GlobalTheory.T.generic_intersection_equals_points_integral}
\formallabel{GlobalTheory.CompanionTheory.Ar.PB.curvatureIntegral_eq_genericIntersection}
\formallabel{PositivityTheory.Xgp.t_deg_generic}
\end{lemma}

\begin{proof}
The generic fibre $P_\Q$ is a projective variety of dimension $D$ over $\Q$
(\lref{GroundTheory.Heights₂.Heights.HilbertSamuel.isProjVarF_map_of_isArithVar}, which states this for every
arithmetic variety). Two field intersection families agree on projective varieties: by induction on the
dimension, (i) fixes the values in dimension $0$, and in dimension $n+1$ \cref{lem:pf:deepstr:stdsec} expresses both
values through values on projective varieties of dimension $n$ (this uniqueness is part of
\lref{GroundTheory.Heights₂.Heights.Targets.ExistsUnique.full}). The intersection numbers of the development
form a field intersection family over $\Q$
(\lref{GroundTheory.Heights₂.Heights.NumbersDVR.IntNumIsFieldIntersection.full}), so $I(P_\Q;\cdot)$ equals them.
Finally, the inputs \ref{G:pf:deepstr:pb} contain the statement of \cref{lem:pf:deepstr:constmetric}(b), and
\lref{GlobalTheory.CompanionTheory.Ar.Giv.degreeGeneric_of_degree} deduces from it (at $p=2$, with the equality of
the degree of the generic fibre and the weighted sum of the degrees of the components of a fibre) that
$\int_Pc_1(L_1)\cdots c_1(L_D)$ is the intersection number of the development of $\deg L_1,\dots,\deg L_D$ on $P_\Q$.
\end{proof}

\begin{theorem}[Nef polarisation]\label{thm:pf:deepstr:nefpol}
Assume the potential-theoretic inputs \ref{G:pf:deepstr:pb}. Let $S$ be an analytic system of dimension $D$ with
index set $\sigma$, $\mathrm{Fl}\subseteq\sigma$ a set of indices whose regularised potentials are pluriharmonic, $y$ a
birational point over a field of characteristic $0$, $\mathrm{Re}$ a realisation of $S$ on $P_y$ with $J=(1)$, and $I$
a field intersection family over $\Q$. Let $n\in\Z^\sigma$ with $n_a\ge0$ for $a\notin\mathrm{Fl}$, and $\alpha=\deg\mathrm{ref}(n)$.
\begin{enumerate}[label=(\alph*)]
\item $\alpha$ is nef on the generic fibre: $I(W;\alpha)\ge0$ for every projective curve $W$ over $\Q$ containing the
generic fibre of $P_y$.
\item If $\deg_S(n)>0$, then for every multihomogeneous $\beta\notin P_y$ of multidegree $d_\beta$ there are $k\ge1$ and
$0\ne e\in\Lambda'_1(y,k\alpha-1-d_\beta)$ (the group of \cref{lem:pf:deepstr:vi}).
\end{enumerate}
\formallabel{GlobalTheory.T.nef_polarization}
\formallabel{GlobalTheory.CompanionTheory.Ar.PB.nefPol}
\end{theorem}

\begin{proof}
(a) Let $W$ be such a curve. It is the generic fibre of an arithmetic curve $W'\supseteq P_y$
(\lref{GlobalTheory.T.arithmetic_variety_generic_fibre_projective}: projective varieties over $\Q$ and
arithmetic varieties correspond under base change and contraction). The geometric degree of $\mathrm{ref}(n)$
on $W'$ is non-negative (\lref{GlobalTheory.T.geometric_degree_nonnegative_psh}: for $n_a\ge0$ off
$\mathrm{Fl}$ the metric of $\mathrm{ref}(n)$ is plurisubharmonic, the indices in $\mathrm{Fl}$ being pluriharmonic, so its
degree on every horizontal curve is $\ge0$), and it equals $I(W;\alpha)$ by
\cref{lem:pf:deepstr:genericfibre}.
(b) The top geometric degree of $\mathrm{ref}(n)$ on $P_y$ is $\deg\mathrm{Re}\cdot\deg_S(n)>0$
(\lref{GlobalTheory.T.realisation_potential_degree_compatibility}, which computes the Chern-form
integrals of reference lines through the realisation), so $I(\text{generic fibre};\alpha^D)>0$ by the same
equality. With (a), the Siu-type statement among \ref{G:pf:deepstr:pb} (a multidegree which is nef on the
generic fibre and has positive top self-intersection there admits, for some $k\ge1$, a nonzero element of
$\Lambda'_1(y,k\alpha-1)$) applies, and \lref{GlobalTheory.T.chart_lattice_shift_nonzero} (which
subtracts any multidegree $d_\beta$ at the cost of increasing $k$) gives $e$.
\end{proof}

\begin{theorem}[Generic nefness on real towers]\label{thm:pf:deepstr:realnef}
Assume the nef polarisation statement and \ref{G:pf:deepstr:pb}. Let $\deepstrII$ be a full intersection system, $T$
a root tower of type $(\Psi,N)$ and $R$ a smooth and flat real tower over $T$. Then for all $B$ and $c$, the model
$M(B,c)$ is nef on the generic fibre at $(\lambda,0)$ for every dominant weight $\lambda$, and at $(0,1)$.
\formallabel{GlobalTheory.T.real_tower_nef_general}
\formallabel{GlobalTheory.CompanionTheory.Ar.Top.rt_nefgen}
\formallabel{RootFieldNumerics.tower_nef_generic_dominant}
\end{theorem}

\begin{proof}
Let $M=M(B,c)$, $\mathrm{Re}$ the realisation of $R$, $I$ the field intersection family over $\Q$ given by the
intersection numbers of the development (\lref{GroundTheory.Heights₂.Heights.NumbersDVR.IntNumIsFieldIntersection.full},
which states that they form such a family), and $\mathrm{Fl}=\{(\tau,N-1)\}$, whose potentials are pluriharmonic by
flatness. Fix an arithmetic curve $W\supseteq P_M$ and put $d(\alpha)=I(W_\Q;\alpha)$; $d$ is additive by
\cref{thm:pf:deepstr:fieldmultilin}, and $\deepstrgdeg(W,L)=d(\deg L)$ for smooth $L$
(\lref{GlobalTheory.CompanionTheory.Ar.Giv.degreeGeneric_of_degree}). By \cref{thm:pf:deepstr:nefpol}(a),
$d(\deg\mathrm{ref}(n))\ge0$ whenever $n_a\ge0$ for $a\notin\mathrm{Fl}$. Taking $n=e_a$ gives $d(\deg\mathrm{ref}(a))\ge0$ for all $a$,
and taking $n=-e_a$ for $a\in\mathrm{Fl}$ gives $d(\deg\mathrm{ref}(a))=0$ there.

For a dominant $\lambda$, $\mathrm{coef}(\lambda,0)_\ast=0$, and with $u_{\tau,k}=-d(a_{(\tau,k)})$ we must show
$\sum_\tau\sum_k\lambda_{\tau,k}u_{\tau,k}\ge0$. The partial sums $U_{\tau,i}=\sum_{k\le i}u_{\tau,k}=d(\deg\mathrm{ref}(\tau,i))$ are $\ge0$, and
$U_{\tau,N-1}=0$. By Abel summation,
$\sum_k\lambda_{\tau,k}u_{\tau,k}=\sum_{i<N-1}(\lambda_{\tau,i}-\lambda_{\tau,i+1})U_{\tau,i}+\lambda_{\tau,N-1}U_{\tau,N-1}\ge0$, since $\lambda$ is
dominant. At $(0,1)$ the sum reduces to $k_A^{-1}\deepstrgdeg(W,A)=k_A^{-1}d(\deg\mathrm{ref}(\ast))\ge0$.
\end{proof}

\subsubsection*{Flag varieties}

\begin{proposition}[Flag degree]\label{prop:pf:deepstr:flagdeg}
For every weight $\lambda$,
\[
\deepstrflagdeg(\lambda)=\deepstrDim!\prod_\tau\prod_{0\le i<k\le N-1}\frac{\lambda_{\tau,i}-\lambda_{\tau,k}}{k-i};
\]
in particular $\deepstrflagdeg(\lambda)>0$ if $\lambda$ is regular dominant and $\deepstrflagdeg(\lambda)\ge0$ if $\lambda$ is dominant.
\formallabel{GlobalTheory.CompanionTheory.Ar.Top.flagDeg_eq_weyl}
\formallabel{Formalism.DegreeFormula.flagDegFormula_statement}
\end{proposition}

\begin{proof}
Write $D=\deepstrDim$. By the multinomial formula, the coefficient of $X^\sigma=\prod_{\tau,i}X_{\tau,i}^{N-1-\sigma_\tau(i)}$ in
$\deepstrlin(\lambda)^D=(-\sum\lambda_{\tau,i}X_{\tau,i})^D$ is $(-1)^DD!\prod_{\tau,i}\lambda_{\tau,i}^{N-1-\sigma_\tau(i)}/\prod_{\tau,i}(N-1-\sigma_\tau(i))!$, and
$\prod_i(N-1-\sigma_\tau(i))!=\prod_{k=0}^{N-1}k!=:s_N$ for every $\tau$. Hence
\[
\deepstrflagdeg(\lambda)=\frac{D!}{s_N^{|\Psi|}}\prod_\tau\sum_{\sigma_\tau\in S_N}\operatorname{sgn}(\sigma_\tau)\prod_i\lambda_{\tau,i}^{N-1-\sigma_\tau(i)}
=\frac{D!}{s_N^{|\Psi|}}\prod_\tau\prod_{i<k}(\lambda_{\tau,i}-\lambda_{\tau,k})
\]
by the Vandermonde determinant, and $\prod_{i<k}(k-i)=s_N$. If $\lambda$ is regular dominant all factors are positive; if
$\lambda$ is dominant they are non-negative.
\end{proof}

\begin{lemma}[Points of given transcendence degree]\label{lem:pf:deepstr:trdegpt}
Let $y=(y_{k,i})$ be a point with coordinates in a field $\Omega$ of characteristic $0$, with at least one block, and
$c$ a chart with $y_{k,c_k}\ne0$ for all $k$. If the subfield of $\Omega$ generated by the ratios $y_{k,i}/y_{k,c_k}$ has
transcendence degree $d$, then $P_y$ is an arithmetic variety of dimension $d+1$.
\formallabel{PositivityTheory.Xgp.fl_arith.arith_of_trdeg}
\end{lemma}

\begin{proof}
The ideal $P_y$ is the kernel of $x_{k,i}\mapsto y_{k,i}T_k$ from $R$ to the domain $\Omega[T_k:k]$, so it is prime; it is
multihomogeneous for the grading by the degrees in the $T_k$, relevant since $x_{k,c_k}\notin P_y$, and $P_y\cap\Z=0$
since $\Omega$ has characteristic $0$. Rescaling $T_k$ by $y_{k,c_k}$ does not change $P_y$, so we may replace $y_{k,i}$ by
$y_{k,i}/y_{k,c_k}$; then the fraction field of $R/P_y$ is $L(T_k:k)$, with $L$ the field of ratios, of transcendence
degree $d+|K|$ (\lref{GroundTheory.Heights₂.Heights.ConeTrdeg.full}, which computes the transcendence degree
of a multicone in this way). As $R/P_y$ is a finitely generated domain containing $\Z$, its Krull dimension is its
transcendence degree plus one (\lref{GroundTheory.Heights₂.Heights.IntDomainDim.full}); so
$\dim R/P_y=d+1+|K|$.
\end{proof}

\begin{lemma}[Generic flags]\label{lem:pf:deepstr:genflag}
Let $F$ be a field, $\Psi$ a finite set, $N\in\N$, and $X_\tau=(X_{\tau,a,b})$, $\tau\in\Psi$, generic $N\times N$ matrices over
$F$. For a block $q=(\tau,k)$, $0\le k\le N-2$, and a $(k+1)$-subset $I\subseteq\{0,\dots,N-1\}$ let $y_{q,I}\in F(X)$ be the
minor of $X_\tau$ on the rows $0,\dots,k$ and the columns $I$, and $y_{q,\mathrm{lead}}=y_{q,\{0,\dots,k\}}$.
\begin{enumerate}[label=(\alph*)]
\item The leading minors $y_{q,\mathrm{lead}}$ are nonzero.
\item The ratios $y_{q,I}/y_{q,\mathrm{lead}}$ generate a field of transcendence degree $|\Psi|\binom N2$ over $F$.
\end{enumerate}
\formallabel{GroundTheory.Heights₂.Heights.LowerBoundFree.GenLeadMinorNeZero.full}
\formallabel{GroundTheory.Heights₂.Heights.LowerBoundFree.FlagRatioTrdeg.full}
\formallabel{GroundTheory.Heights₂.Heights.LowerBoundFree.FlagRatioTrdeg.proof}
\end{lemma}

\begin{proof}
(a) Specialising every $X_\tau$ to the identity matrix maps a leading minor to $1$, so it is a nonzero polynomial.
(b) By the same specialisation all leading principal minors of $X_\tau$, of orders $1,\dots,N$, are nonzero, so Gaussian
elimination without row exchanges over $F(X)$ gives $X_\tau=\Lambda_\tau U_\tau$ with $\Lambda_\tau$ lower triangular and invertible
and $U_\tau$ upper unitriangular, with entries in $F(X)$. The entries of the $\Lambda_\tau$ on and below the diagonal and of
the $U_\tau$ above it, $|\Psi|N^2$ elements in all, generate $F(X)$ (as $X_\tau=\Lambda_\tau U_\tau$), whose transcendence degree
is $|\Psi|N^2$; hence they are algebraically independent. For $q=(\tau,k)$, the rows $0,\dots,k$ of $X_\tau$ are the upper
left $(k+1)\times(k+1)$ block $\Lambda_{\tau,k}$ of $\Lambda_\tau$ times the rows $0,\dots,k$ of $U_\tau$. Writing $U_\tau(k;I)$ for the
submatrix of $U_\tau$ on the rows $0,\dots,k$ and the columns $I$, we get $y_{q,I}=\det\Lambda_{\tau,k}\cdot\det U_\tau(k;I)$, and
since $U_\tau(k;\{0,\dots,k\})$ is unitriangular, $y_{q,I}/y_{q,\mathrm{lead}}=\det U_\tau(k;I)$. So all ratios lie in the field
generated by the entries $U_{\tau,a,b}$, $a<b$. Conversely, $U_{\tau,a,b}$ is the ratio for $q=(\tau,a)$ and
$I=\{0,\dots,a-1,b\}$ (here $a\le N-2$ as $b\le N-1$): the matrix $U_\tau(a;I)$ is unitriangular in its first $a$ columns
and its last column is $(U_{\tau,0,b},\dots,U_{\tau,a,b})$. Hence the ratios generate the field generated by the
$|\Psi|\binom N2$ algebraically independent elements $U_{\tau,a,b}$, $a<b$.
\end{proof}

\begin{proposition}[Dimension drop at the boundary of the flag variety]\label{prop:pf:deepstr:flagdrop}
Let $\Psi\ne\varnothing$ and $N\ge2$. For $\tau\in\Psi$ let $X_\tau=(X_{\tau,a,b})$ be a generic $N\times N$ matrix over $\Q$; for a
\emph{block} $q=(\tau,k)$, $0\le k\le N-2$, and a $(k+1)$-subset $I\subseteq\{0,\dots,N-1\}$ let $y_{q,I}\in\Q(X)$ be the minor of
$X_\tau$ on the rows $0,\dots,k$ and the columns $I$ (the Pl\"ucker coordinates of the generic flag), and let $P_{\mathrm{Fl}}$ be
the ideal of this point. Then $P_{\mathrm{Fl}}$ is an arithmetic variety of dimension $\deepstrDim+1$, and for every block
$q_0$ such that another block exists, the image of $P_{\mathrm{Fl}}$ under the coordinate projection forgetting the block
$q_0$ is an arithmetic variety of dimension $<\deepstrDim+1$.
\formallabel{GlobalTheory.T.flag_variety_boundary_dimension_drop}
\formallabel{GlobalTheory.CompanionTheory.Ar.PB.flag_isArithVar_and_image}
\formallabel{PositivityTheory.Xgp.t_fl_arith}
\formallabel{PositivityTheory.Xgp.fl_arith.trdeg_lt}
\formallabel{GroundTheory.Heights₂.Heights.LowerBoundFree.FlagArithVar.proof}
\end{proposition}

\begin{proof}
Note $\deepstrDim=|\Psi|\binom N2$. By \cref{lem:pf:deepstr:genflag} with $F=\Q$, the leading minors $y_{q,\mathrm{lead}}$
(columns $0,\dots,k$) are nonzero and the field $L$ generated by all ratios $y_{q,I}/y_{q,\mathrm{lead}}$ has transcendence
degree $\deepstrDim$; so the first assertion follows from \cref{lem:pf:deepstr:trdegpt}. For the second, the ideal of the
image is the ideal of the point obtained by forgetting the coordinates of $q_0$. Let $L'\subseteq L$ be generated by the
ratios of the blocks $q\ne q_0$, and $t=y_{q_0,\mathrm{next}}/y_{q_0,\mathrm{lead}}\in L$, where $\mathrm{next}$ is the set of columns
$\{0,\dots,k-1,k+1\}$ for $q_0=(\tau_0,k)$. For $c\in\Q$ let $\sigma_c$ be the automorphism of $\Q(X)$ adding $c$ times the
row $k+1$ of $X_{\tau_0}$ to its row $k$. It fixes every minor of a block $q\ne q_0$: for $q=(\tau,j)$ with $\tau\ne\tau_0$ or
$j<k$ the rows involved are unchanged, and for $j>k$ both rows $k$ and $k+1$ are among the rows $0,\dots,j$, so the
minor is unchanged. Hence $\sigma_c$ fixes $L'$, while an evaluation of the variables shows that the values
$\sigma_n(t)$, $n\in\N$, are pairwise distinct. Therefore $t$ has infinitely many conjugates over $L'$, so it is
transcendental over $L'$, and $\operatorname{trdeg}L'=d<\deepstrDim$ because $L'(t)\subseteq L$. Finally, the point
obtained by forgetting $q_0$ has nonzero leading coordinates and its ratios generate $L'$, so by \cref{lem:pf:deepstr:trdegpt}
the image is an arithmetic variety of dimension $d+1<\deepstrDim+1$.
\end{proof}

\subsubsection*{An integral section attached to an origin}

\begin{proposition}[Weight section]\label{prop:pf:deepstr:weightsection}
Let $\mathrm{Or}$ be a root origin of a root model $M$ over a model $Y$, relative to a finite set $S$ of primes and
fixed data $F_x$ with integers $k_A,q'\ge1$, exponents $n_r$, $r\in\Psi\times\{0,\dots,N-1\}$, and a nonzero element $s_E$, and
let $\nu$ be a weight with $n_r=k_Aq'\nu_r$ for all $r$. Put $n_0(r)=-n_r$ for $r\ne\ast$ and $n_0(\ast)=-q'$, so that
$n_0=k_Aq'\,\mathrm{coef}(\nu,-1)$. Then there is an integer $N_0\ge1$ such that
$\varepsilon=N_0\,g(s_E)\prod_re_r^{n_r}\in\Omega_M$ is nonzero and lies in $\Lambda_1(y,\deg L_1(n_0))$, where $g$ and $e_r$ are the field
embedding and the root elements of the origin.
\formallabel{GlobalTheory.T.root_origin_weight_identity}
\formallabel{GlobalTheory.CompanionTheory.Ar.Small.OriginEps.section_on_model}
\end{proposition}

\begin{proof}
The identity $n_0=k_Aq'\,\mathrm{coef}(\nu,-1)$ holds by the definitions ($k_A$ of $M$ equals that of $F_x$). The element
$s_0=g(s_E)\prod_re_r^{n_r}$ is nonzero since $s_E\ne0$ and the $e_r$ are nonzero. Let $c'$ be a chart with all
$y_{k,c'_k}\ne0$ and $a=\deg L_1(n_0)=-\sum_rn_ra_r-q'\deg A$. A direct computation, using that the line $A$ has degree $1$ in the
block $k_0$ of the coordinates of $F_x$ and $0$ in the others, gives
\[
s_0\prod_ky_{k,c'_k}^{-a_k}=\frac{g(s_E)\,g(y_f)^{q'}}{\prod_r\bigl(\prod_ky_{k,c'_k}^{(a_r)_k}\cdot e_r\bigr)^{-n_r}},
\]
with $y_f$ the coordinate of $F_x$ in the chart. At every point of the chart outside the closed fibres this
quotient is regular, by the effectivity condition of the root origin. Elements regular at all such points
lie in the $\Q$-algebra generated by the chart ratios
(\lref{GlobalTheory.CompanionTheory.Ar.Small.OriginEps.mem_adjoin_rat_of_regular}), and the latter is contained
in the algebra generated by the ratios $y_p/y_{p,c'}$, so $s_0\prod_ky_{k,c'_k}^{-a_k}$ is integral over that algebra.
By \lref{GlobalTheory.T.integral_section_lattice_multiple} (an element integral over the rational chart
algebras on all charts has a positive integer multiple in $\Lambda_1(y,a)$) there is $N_0\ge1$ with $N_0s_0\in\Lambda_1(y,a)$.
\end{proof}

\subsubsection*{Hodge--Tate points and charts}

\begin{definition}[Hodge--Tate points]\label{def:pf:deepstr:pointht}
The \emph{Hodge--Tate point statement} $\mathrm{(HTP)}$ is: let $p$ be a prime, $C$ an algebraically closed field of characteristic $0$
with a valuation $v$, $v(p)<1$, whose valuation ring $\OO$ is a $p$-adically complete domain; $E$ a number field with
an involution $c$; $G$ a commutative semi-abelian group scheme over a locally noetherian scheme $S$, and
$y\in S(\OO)$; $\Lambda$ a lattice with an $\OO_E$-action and an alternating form $\langle\ ,\ \rangle$ which is polarised at $p$ with
respect to a basis $u_1,\dots,u_N$; $X$ a tower of level structures identifying the $p$-power torsion of $G_y$ with
$\Lambda$, compatible with an $\OO_E$-action on $G$ and on a Lie model $\mathcal L$ of $G$ at $y$ with
$2\rank_{\OO}\mathcal L=\rank_\Z\Lambda$; and $(\zeta_n)$ a compatible system of primitive $p^n$-th roots of unity for which the
Weil pairing of $X$ is regular. Then for every $\kappa\in\OO$ with $\kappa\ne0$ and $v(\kappa)^{p-1}\le v(p)$ there is a
\emph{Hodge--Tate point} of level $\kappa$ for $(C\otimes\Lambda,C\otimes\mathcal L^\vee)$ attached to $X$: its lattice has basis
$1\otimes u_i$, its form is $\langle\ ,\ \rangle$, its Hodge lattice is $\OO\otimes\mathcal L^\vee$, its $\OO_E$-actions are the given ones
(twisted by $c$ on $\mathcal L^\vee$), and its Hodge--Tate map sends $1\otimes x$ into $1\otimes w+p^n\cdot(\text{Hodge lattice})$ whenever
$w$ is the derivative at level $n$ of the character $\zeta_n^{\langle x,\cdot\rangle}$ of the level group.
\formallabel{GlobalTheory.CompanionTheory.Ar.HT.PointHT}
\end{definition}

\begin{lemma}[Integers of complete algebraically closed fields]\label{lem:pf:deepstr:integers}
Let $p$ be a prime, $C$ an algebraically closed field of characteristic $0$ with a valuation $v\colon C\to\R_{\ge0}$
such that $v(p)<1$, and $\OO=\{v\le1\}$ its ring of integers, assumed $p$-adically complete. Then $\OO$ is a valuation
ring with fraction field $C$, hence integrally closed with algebraically closed fraction field; it is a Henselian
local ring and a $\Z_p$-algebra.
\formallabel{GlobalTheory.T.complete_integers_properties}
\formallabel{GlobalTheory.CompanionTheory.Ar.HT.padicComplete_integers_structure}
\formallabel{PositivityTheory.HtPt.Twhopf.padic_alg}
\end{lemma}

\begin{proof}
For $x\in C^\times$ we have $v(x)\le1$ or $v(x^{-1})\le1$, so $\OO$ is a valuation ring with fraction field $C$; valuation
rings are local and integrally closed. Let $f\in\OO[X]$ be monic and $a_0\in\OO$ with $f(a_0)$ in the maximal ideal
$\mathfrak m$ and $f'(a_0)\notin\mathfrak m$. Over $C$, $f=\prod_i(X-\alpha_i)$, and the $\alpha_i$ are integral over $\OO$, hence in $\OO$.
Modulo $\mathfrak m$, $\bar f=\prod_i(X-\bar\alpha_i)$ has the simple root $\bar a_0$, so exactly one $\alpha_i$ reduces to $\bar a_0$, and it is a
root of $f$ lifting $a_0$; thus $\OO$ is Henselian. Finally the maps $\Z\to\OO/p^n\OO$ factor through $\Z/p^n\Z$ and are
compatible, so they define $\Z_p=\varprojlim\Z/p^n\Z\to\varprojlim\OO/p^n\OO=\OO$, the last equality because $\OO$ is $p$-adically
complete (\lref{GlobalTheory.ChF.padicInt_algebra_of_isAdicComplete_and_matrix_unique}, which constructs
this algebra structure).
\end{proof}

\begin{theorem}[Hodge--Tate points]\label{thm:pf:deepstr:pointht}
The sharp Hodge--Tate statement \ref{G:pf:deepstr:ht} implies $\mathrm{(HTP)}$.
\formallabel{GlobalTheory.T.pointHT_of_fargues_hodge_tate}
\formallabel{GlobalTheory.CompanionTheory.Ar.HT.Fargues.pointHT}
\end{theorem}

\begin{proof}
Let the data be as in \cref{def:pf:deepstr:pointht}. By \cref{lem:pf:deepstr:integers}, $\OO$ is a Henselian local
$\Z_p$-algebra with algebraically closed fraction field $C$; since $G$ is smooth and $\OO$ is local, the Lie model
$\mathcal L$ is free of finite rank. From $X$ one forms the tower of $p$-power torsion of $G$ at $y$
over $\OO$ (\lref{GlobalTheory.CompanionTheory.Ar.HT.Fargues.ordTower}). The $\OO_E$-action is compatible with
this tower (\lref{GlobalTheory.T.semiabelian_adic_torsion_four}, which transports homomorphisms of $G$
inducing the given maps on $\Lambda$ and $\mathcal L$ to the tower), and the pairing extends with values in powers of
$\zeta_n$ (\lref{GlobalTheory.T.exists_prodSpan_torsion_pairing_eq_zeta_pow}). The existence of the
Hodge--Tate point is then \lref{GlobalTheory.T.polarised_divisible_group_exists}, which constructs it
from a polarised $p$-divisible group with these structures, using the sharp Hodge--Tate statement for the
cokernel of the Hodge--Tate map.
\end{proof}

\begin{theorem}[Hodge--Tate charts]\label{thm:pf:deepstr:htcharts}
Let $(E,H,D_w,l,X^o,c)$ be as in \cref{sec:pf:deepstr:conv}, $Y$ a model of $X^o$ over a ring mapping to $F_c$, $S$ a
finite set of primes, $y$ a point in coordinates with field $\Omega$, $\mathfrak M$ a Hodge component of $Y$ at $S$ in these
coordinates, $p\in S$ with $p+q>0$, and $o\subseteq\Omega$ the ring of fractions $a/b$ with $p\nmid b$. Then there is a family
$(\mathrm{Dc}_s)$ of Hodge--Tate charts of level $p^{R(p)}$ and rank $p+q$, one over the chart ring of each standard chart
$s$ of the multicone (with values in the dual of the Lie module of $\mathfrak M$), which forms a Hodge--Tate family for
$\mathfrak M$ at $p$, and, if $\mathrm{(HTP)}$ holds, each $\mathrm{Dc}_s$ has test points at every element of its chart ring
which is a $\C_p$-element.
\formallabel{GlobalTheory.T.hodge_tate_free_flag_lattice_lower}
\formallabel{GlobalTheory.CompanionTheory.Ar.HT.hodgeTate_family_exists}
\formallabel{GlobalTheory.CompanionTheory.HtLink.htat_link_of}
\end{theorem}

\begin{proof}
Generic properties of the Hodge component (among them properness of the semi-abelian scheme over its
base, Tate isomorphisms, the rank of the Lie module and a description of each standard chart) are
\lref{PositivityTheory.HtPt.Hcptgen.hcpt_generic}; the Hodge component also has full $p^{R(p)}$-torsion.
They provide a level parameter (\lref{PositivityTheory.HtPt.Htat.exists_isParam}); the lattice of the datum
has a basis polarised at $p$ (\lref{PositivityTheory.HtPt.Htat.pol_lattice}), which gives a level pairing
(\lref{PositivityTheory.HtPt.Htat.exists_levPairing}). For each standard chart $s$,
\lref{GlobalTheory.CompanionTheory.HtLink.exists_chart} constructs a Hodge--Tate chart over the chart ring
from these data and the existence of Hodge--Tate charts (\lref{GlobalTheory.T.HTChartExists_holds}); its
Hodge lattice spans the Hodge lattices of $\mathfrak M$ at the points of the chart over $\Z_p$, and its Hodge--Tate map
describes the Hodge--Tate condition at these points. The charts share the same level pairing, and the compatibility conditions of a
Hodge--Tate family follow from these descriptions, the chart rings being contained in the local rings of the
points (\lref{PositivityTheory.HtPt.Htat.chart_le_regAt}). The test points come from
\lref{GlobalTheory.CompanionTheory.Ar.HT.hasTestPoints_of_tate}, which deduces them from $\mathrm{(HTP)}$.
\end{proof}

% ---------------------------------------------------------------------------
\subsection{Remarks on the formal statements}\label{sec:pf:deepstr:remarks}

\begin{remark}[Use of the hypotheses]\label{rem:pf:deepstr:hyps}
(1) Within the proof of \cref{pf:deepstr:main}, asymptotic Riemann--Roch \ref{H:pf:deepstr:nef} is used only
through \cref{prop:pf:deepstr:geominputs}\ref{gi:topdeg}, the identification of the top degree of a component
with the volume of the corresponding root model. Zarhin's trick \ref{H:pf:deepstr:zar} and the existence of
quotients by finite flat subgroups enter only through the good tower volume theorem.

(2) Of the hypothesis $m^2\le k$ of $\sfname{DeepStR}$ only the consequence $k\ge2m-1$ is used: at level $j\ge1$
the condition $3q+1\le p+q$ reads $2m-3j+2\le k$. The bound $15$ on the number of prime factors of the level
modulus enters only through the inequality $\#\mathcal S\le n+19\le128n$ in the proof of
\cref{thm:pf:deepstr:ofvolume}.

(3) The constant of $\sfname{DeepStR}$ is $C_T+120$, where $C_T$ is the constant of $\sfname{VolumeTheoremX}$;
the summand $120=2\cdot60$ comes from \cref{lem:pf:deepstr:costbound}.
\end{remark}

\begin{remark}[Peculiarities]\label{rem:pf:deepstr:peculiar}
(1) The volume inequality $\mathrm{(VI)}$ of \cref{def:pf:deepstr:bppf} is called a birational pull-back and
push-forward formula in the development; its content is the inequality between arithmetic self-intersection
and arithmetic volume stated there. It is used in this section only through \cref{lem:pf:deepstr:shift}(b),
and in the development it is a consequence of $\sfname{MoriwakiBigStatement}$ (\cref{lem:pf:deepstr:vi}).

(2) The third component of the Zhang-type provisions \ref{G:pf:deepstr:zhang} is the true proposition and
plays no role. The formal conclusion of \cref{thm:pf:deepstr:slopebound} states the Lie slope bound twice.

(3) In \cref{lem:pf:deepstr:counting}(b) the number of elements of an infinite set is taken to be $0$, as in
the development; the inequality then holds trivially. Divisions by $0$ have the value $0$; in
\cref{lem:pf:deepstr:towerseq} this convention makes both sides of the second limit vanish when the geometric
degree vanishes, which happens only for finitely many $n$.

(4) The limit functional $\mathrm{PP}$ of a root model is defined as a limit as $t\to0^+$; when the limit
does not exist its value is an unspecified real number. In this section $\mathrm{PP}$ is only used through
the limit conditions $\mathrm{(LC)}$ of the good tower volume theorem, which include the convergence of the
relevant limits.

(5) $\sfname{DeepStR}$ is a statement about the volume slopes $\deepstrvsl(Z,u)$ of \eqref{eq:pf:deepstr:compvol},
which are defined with the counting volumes at the step $a=8$ and the valuation families $\mathcal V_B$ with
$B=\deepstrbnd(u,g)$; the error term involves $\log\deepstrbnd(u,g)$ through this choice of $B$.

(6) Two proofs in \cref{sec:pf:deepstr:further} take a shorter route than the development. The development
obtains \cref{lem:pf:deepstr:constmetric}(b) from the conservation of degrees in the fibres rather than from part
(a) and the recursion, and it obtains \cref{lem:pf:deepstr:genflag}(b) from the algebraic independence of the
minors on the rows $0,\dots,a$ and the columns $\{0,\dots,a-1,b\}$, $a\le b$, rather than from the factorisation
$X_\tau=\Lambda_\tau U_\tau$; the two routes prove the same statements.
\end{remark}

